\begin{proof}[Proof of \cref{obj:prop:identification}]
\begin{enumerate}
\item \emph{Completion.}
Fix \(\mathbb P\in\mathcal D_{d,\epsilon}^{\mathrm{obs}}\), and use its cell quantities from \cref{obj:def:observed-class}. Define the Bernoulli probability weights by
\[
\operatorname{Bern}(\vartheta;y)
=
\begin{cases}
1-\vartheta,&y=0,\\
\vartheta,&y=1,
\end{cases}
\qquad
\vartheta\in[0,1],\quad y\in\{0,1\}.
\]
For \(x\in[d]\) and \(a,y,y_0,y_1\in\{0,1\}\), define the full-data atom probabilities by
\[
\begin{aligned}
&P(X=x,A=a,Y=y,Y(0)=y_0,Y(1)=y_1)\\
&\qquad=
\begin{cases}
q_{ay,x}\operatorname{Bern}(\mu_{1-a,x};y_{1-a}),&y=y_a,\\
0,&y\ne y_a,
\end{cases}
\end{aligned}
\]
where
\[
y_a=
\begin{cases}
y_0,&a=0,\\
y_1,&a=1,
\end{cases}
\qquad
y_{1-a}=
\begin{cases}
y_1,&a=0,\\
y_0,&a=1.
\end{cases}
\]
The bounds \(q_{ay,x}\ge0\) and \(\mu_{ax}\in[0,1]\) from \cref{obj:def:observed-class} make these masses nonnegative. Summing first over the potential-outcome coordinates gives
\[
\begin{aligned}
&\sum_{x\in[d]}\sum_{a,y,y_0,y_1\in\{0,1\}}
P(X=x,A=a,Y=y,Y(0)=y_0,Y(1)=y_1)\\
&\quad=
\sum_{x\in[d]}\sum_{a,y\in\{0,1\}}
q_{ay,x}
\bigl\{
\operatorname{Bern}(\mu_{1-a,x};0)
+\operatorname{Bern}(\mu_{1-a,x};1)
\bigr\}\\
&\quad=\sum_{x\in[d]}\sum_{a,y\in\{0,1\}}q_{ay,x}=1.
\end{aligned}
\]
Thus the construction defines a probability law. Its atoms with \(y\ne y_a\) have zero mass, so
\[
P\{Y=Y(A)\}=1,
\]
which is \cref{obj:ass:consistency}.
% lean: CausalSmith.Stat.OptvalueVanishingoverlapRate.completionFullMass_nonneg_sum
% lean: CausalSmith.Stat.OptvalueVanishingoverlapRate.bernoulliCompletion_consistency

Summing over the observed outcome, and then using the observed-atom factorization supplied by \cref{obj:def:observed-class}, yields
\[
\begin{aligned}
&P(X=x,A=a,Y(0)=y_0,Y(1)=y_1)\\
&\quad=q_{a\,y_a,x}\operatorname{Bern}(\mu_{1-a,x};y_{1-a})\\
&\quad=p_x\operatorname{Bern}(\pi_x;a)
\operatorname{Bern}(\mu_{0x};y_0)
\operatorname{Bern}(\mu_{1x};y_1).
\end{aligned}
\]
Consequently, summing these factors over the remaining binary coordinates gives, for \(a,a',y\in\{0,1\}\),
\[
\begin{aligned}
P(X=x)&=p_x,\\
P(X=x,A=a')&=p_x\operatorname{Bern}(\pi_x;a'),\\
P(X=x,Y(a)=y)&=p_x\operatorname{Bern}(\mu_{ax};y),\\
P(X=x,A=a',Y(a)=y)
&=p_x\operatorname{Bern}(\pi_x;a')
\operatorname{Bern}(\mu_{ax};y).
\end{aligned}
\]
Multiplication therefore gives
\[
\begin{aligned}
&P(X=x,A=a',Y(a)=y)\,P(X=x)\\
&\qquad=P(X=x,A=a')\,P(X=x,Y(a)=y).
\end{aligned}
\]
On positive-mass cells, division by \(p_x^2\) proves \cref{obj:ass:conditional-exchangeability}; on null cells both products are zero.
% lean: CausalSmith.Stat.OptvalueVanishingoverlapRate.bernoulliCompletion_poAtom_factorization
% lean: CausalSmith.Stat.OptvalueVanishingoverlapRate.bernoulliCompletion_exchangeability

For each observed atom, the construction also gives
\[
\begin{aligned}
P(X=x,A=a,Y=y)
&=q_{ay,x}
\bigl\{
\operatorname{Bern}(\mu_{1-a,x};0)
+\operatorname{Bern}(\mu_{1-a,x};1)
\bigr\}\\
&=q_{ay,x}.
\end{aligned}
\]
Hence
\[
\operatorname{Obs}(P)=\mathbb P.
\]
The observed marginal inherits the overlap restriction of \(\mathbb P\). Together with consistency and conditional exchangeability, this puts \(P\) in the class of \cref{obj:def:causal-class}, proving the completion assertion.
% lean: CausalSmith.Stat.OptvalueVanishingoverlapRate.bernoulliCompletion_observedMarginal

\item \emph{Value identification.}
Fix \(P\in\mathcal P_{d,\epsilon}\), and use \(q_{ay,x},p_x,\pi_x,\mu_{ax}\) for the quantities of its observed marginal defined in \cref{obj:def:observed-class}. Marginalization gives
\[
P(X=x)
=\sum_{a,y\in\{0,1\}}
\operatorname{Obs}(P)(X=x,A=a,Y=y)
=p_x.
\]
We compare the two value sums cell by cell. If \(p_x=0\), their cell contributions are both zero:
\[
p_x\max_{a\in\{0,1\}}\mathbb E_P[Y(a)\mid X=x]
=
p_x\max_{a\in\{0,1\}}\mu_{ax}
=0.
\]
Here the conditional potential-outcome means use the stated zero convention.
% lean: CausalSmith.Stat.OptvalueVanishingoverlapRate.observedMarginal_cellMass

Suppose now that \(p_x>0\). The lower overlap bound in \cref{obj:ass:shrinking-overlap} gives
\[
q_{10,x}+q_{11,x}=p_x\pi_x\ge\epsilon p_x>0.
\]
The decomposition of the cell mass and the upper overlap bound give
\[
\begin{aligned}
p_x&=(q_{00,x}+q_{01,x})+(q_{10,x}+q_{11,x}),\\
q_{10,x}+q_{11,x}&\le(1-\epsilon)p_x,
\end{aligned}
\]
and hence
\[
q_{00,x}+q_{01,x}
=p_x-(q_{10,x}+q_{11,x})
\ge\epsilon p_x>0.
\]
Thus both arm masses are positive.

For either \(a\in\{0,1\}\), \cref{obj:ass:conditional-exchangeability} supplies
\[
P(X=x,A=a,Y(a)=1)\,p_x
=
P(X=x,A=a)\,P(X=x,Y(a)=1).
\]
Consistency, as specified in \cref{obj:ass:consistency}, identifies the selected atoms:
\[
P(X=x,A=a,Y(a)=y)=q_{ay,x},
\qquad y\in\{0,1\}.
\]
Summing this identity over \(y\), and substituting its outcome-one instance into the exchangeability identity, gives
\[
P(X=x,A=a)=q_{a0,x}+q_{a1,x},
\qquad
q_{a1,x}p_x
=(q_{a0,x}+q_{a1,x})P(X=x,Y(a)=1).
\]
Division by the positive cell and arm masses now yields
\[
\mathbb E_P[Y(a)\mid X=x]
=
\frac{P(X=x,Y(a)=1)}{p_x}
=
\frac{q_{a1,x}}{q_{a0,x}+q_{a1,x}}
=\mu_{ax}.
\]
Combining the positive- and zero-mass cases and summing over \(x\) proves
\[
\begin{aligned}
V^\star(P)
&=\sum_{x\in[d]}P(X=x)
\max_{a\in\{0,1\}}\mathbb E_P[Y(a)\mid X=x]\\
&=\sum_{x\in[d]}p_x\max_{a\in\{0,1\}}\mu_{ax}
=\Psi\{\operatorname{Obs}(P)\},
\end{aligned}
\]
where the last equality is the definition in \cref{obj:def:observed-value}.
% lean: CausalSmith.Stat.OptvalueVanishingoverlapRate.poRegression_eq_outcomeMean_of_pos
% lean: CausalSmith.Stat.OptvalueVanishingoverlapRate.identification

\item \emph{The armwise extension on observed cells.}
Fix \(\mathbb P\in\mathcal D_{d,\epsilon}^{\mathrm{obs}}\) and \(x\in[d]\). The cell vector and the quantities in \cref{obj:def:armwise-extension} satisfy
\[
s(q_x)=p_x,\qquad
s_a(q_x)=q_{a0,x}+q_{a1,x},\qquad
(q_x)_{a1}=q_{a1,x},
\qquad a\in\{0,1\}.
\]
Every coordinate is nonnegative, so \(p_x\ge0\). If \(p_x=0\), then
\[
0\le q_{ay,x}\le p_x=0
\qquad(a,y\in\{0,1\}),
\]
and therefore \(q_x=0\). The zero branch of \cref{obj:def:armwise-extension} gives
\[
F_\epsilon(q_x)=0
=p_x\max_{a\in\{0,1\}}\mu_{ax}.
\]

Suppose instead that \(p_x>0\). The lower overlap bound gives
\[
\epsilon p_x\le q_{10,x}+q_{11,x}=s_1(q_x).
\]
For the other arm, the upper overlap bound and the cell-mass decomposition give
\[
s_1(q_x)\le(1-\epsilon)p_x,
\qquad
s_0(q_x)=p_x-s_1(q_x)\ge\epsilon p_x.
\]
Thus, for each \(a\in\{0,1\}\),
\[
s_a(q_x)\ge\epsilon p_x>0,
\qquad
D_a(q_x)=\max\{s_a(q_x),\epsilon p_x\}=s_a(q_x).
\]
Each armwise term consequently reduces to
\[
\frac{s(q_x)(q_x)_{a1}}{D_a(q_x)}
=
\frac{p_xq_{a1,x}}{q_{a0,x}+q_{a1,x}}
=
p_x\mu_{ax}.
\]
Since \(s(q_x)=p_x>0\), the vector \(q_x\) is nonzero. Using the nonzero branch of \cref{obj:def:armwise-extension} and factoring the nonnegative \(p_x\) out of the maximum gives
\[
F_\epsilon(q_x)
=
\max_{a\in\{0,1\}}p_x\mu_{ax}
=
p_x\max_{a\in\{0,1\}}\mu_{ax}.
\]
This proves the displayed armwise identity, including its null-cell case.
% lean: CausalSmith.Stat.OptvalueVanishingoverlapRate.armwiseExtension_cellVector
\end{enumerate}
\end{proof}

\begin{proof}[Proof of \cref{obj:thm:matched-frontier}]
\begin{enumerate}
\item Choose the universal constants \(c_\star,C_\star,H_0,\kappa,D_0\) supplied by the risk guarantees in \cref{obj:thm:weighted-separation-and-frontier}. They satisfy
\[
0<c_\star\le C_\star,\qquad H_0>0,\qquad \kappa>0,\qquad D_0\ge2.
\]
Fix integers \(n\ge1\), \(d\ge2\), and an overlap parameter \(0<\epsilon\le1/2\). Define
\[
L_d=\log(ed),
\qquad
\rho_{n,d,\epsilon}
=\min\left\{1,\frac{d}{n\epsilon L_d}\right\}.
\]
For the estimator of \cref{obj:def:armwise-estimator} with the chosen calibrations, \cref{obj:thm:weighted-separation-and-frontier} gives
\[
c_\star\rho_{n,d,\epsilon}
\le \mathfrak R_{n,d,\epsilon}
\le
\sup_{\mathbb P\in\mathcal D_{d,\epsilon}^{\mathrm{obs}}}
E_{\mathbb P}\!\left[
\bigl(\widehat V_{n,d,\epsilon}^{\mathrm{AF}}-\Psi(\mathbb P)\bigr)^2
\right]
\le C_\star\rho_{n,d,\epsilon},
\]
and
\[
c_\star\rho_{n,d,\epsilon}
\le
\inf_{\widehat V}
\sup_{P\in\mathcal P_{d,\epsilon}}
E_{\operatorname{Obs}(P)}\!\left[
\bigl(\widehat V-V^\star(P)\bigr)^2
\right]
\le C_\star\rho_{n,d,\epsilon}.
\]
The infimum is over measurable real-valued estimators based on the observed sample. Thus the observed chain and the two causal minimax bounds are already established with the same universal constants.
% lean: CausalSmith.Stat.OptvalueVanishingoverlapRate.weighted_separation_and_frontier

\item We next establish equality of the causal and observed worst-case risks for each estimator. Fix any measurable function
\[
\widehat V:
\bigl([d]\times\{0,1\}\times\{0,1\}\bigr)^n
\longrightarrow\mathbb R.
\]
By \cref{obj:prop:identification}, every \(P\in\mathcal P_{d,\epsilon}\) satisfies
\[
V^\star(P)=\Psi\{\operatorname{Obs}(P)\}.
\]
Consequently,
\[
E_{\operatorname{Obs}(P)}\!\left[
\bigl(\widehat V-V^\star(P)\bigr)^2
\right]
=
E_{\operatorname{Obs}(P)}\!\left[
\bigl(\widehat V-\Psi\{\operatorname{Obs}(P)\}\bigr)^2
\right].
\]
Here both expectations use the independent sample of size \(n\) from the observed marginal, as in \cref{obj:ass:iid-sampling}.
% lean: CausalSmith.Stat.OptvalueVanishingoverlapRate.causalRisk_eq_observedRisk

Every causal law has an observed marginal in \(\mathcal D_{d,\epsilon}^{\mathrm{obs}}\) by \cref{obj:def:causal-class}. Conversely, \cref{obj:prop:identification} supplies a causal completion for every law in that observed class. The displayed risk identity therefore gives both inclusions in the equality of attained risk sets:
\[
\begin{aligned}
&\left\{
E_{\operatorname{Obs}(P)}\!\left[
\bigl(\widehat V-V^\star(P)\bigr)^2
\right]:
P\in\mathcal P_{d,\epsilon}
\right\}
\\
&\qquad=
\left\{
E_{\mathbb P}\!\left[
\bigl(\widehat V-\Psi(\mathbb P)\bigr)^2
\right]:
\mathbb P\in\mathcal D_{d,\epsilon}^{\mathrm{obs}}
\right\}.
\end{aligned}
\]
% lean: CausalSmith.Stat.OptvalueVanishingoverlapRate.causalRisk_range_eq_observedRisk_range
Taking suprema yields
\[
\sup_{P\in\mathcal P_{d,\epsilon}}
E_{\operatorname{Obs}(P)}\!\left[
\bigl(\widehat V-V^\star(P)\bigr)^2
\right]
=
\sup_{\mathbb P\in\mathcal D_{d,\epsilon}^{\mathrm{obs}}}
E_{\mathbb P}\!\left[
\bigl(\widehat V-\Psi(\mathbb P)\bigr)^2
\right].
\]
% lean: CausalSmith.Stat.OptvalueVanishingoverlapRate.causalWorstRisk_eq_observedWorstRisk

\item Taking the infimum of this identity over the common estimator class gives
\[
\inf_{\widehat V}
\sup_{P\in\mathcal P_{d,\epsilon}}
E_{\operatorname{Obs}(P)}\!\left[
\bigl(\widehat V-V^\star(P)\bigr)^2
\right]
=\mathfrak R_{n,d,\epsilon}.
\]
% lean: CausalSmith.Stat.OptvalueVanishingoverlapRate.causalMinimaxRisk_eq_minimaxRisk
Applying the same identity to the prescribed armwise estimator gives
\[
\begin{aligned}
&\sup_{P\in\mathcal P_{d,\epsilon}}
E_{\operatorname{Obs}(P)}\!\left[
\bigl(\widehat V_{n,d,\epsilon}^{\mathrm{AF}}-V^\star(P)\bigr)^2
\right]
\\
&\qquad=
\sup_{\mathbb P\in\mathcal D_{d,\epsilon}^{\mathrm{obs}}}
E_{\mathbb P}\!\left[
\bigl(\widehat V_{n,d,\epsilon}^{\mathrm{AF}}-\Psi(\mathbb P)\bigr)^2
\right].
\end{aligned}
\]
% lean: CausalSmith.Stat.OptvalueVanishingoverlapRate.causalArmwiseWorstRisk_eq_armwiseWorstRisk
Substituting these two equalities into the middle observed inequality from the first step, and then substituting the second equality into its upper inequality, yields
\[
\begin{aligned}
\inf_{\widehat V}
\sup_{P\in\mathcal P_{d,\epsilon}}
E_{\operatorname{Obs}(P)}\!\left[
\bigl(\widehat V-V^\star(P)\bigr)^2
\right]
&\le
\sup_{P\in\mathcal P_{d,\epsilon}}
E_{\operatorname{Obs}(P)}\!\left[
\bigl(\widehat V_{n,d,\epsilon}^{\mathrm{AF}}-V^\star(P)\bigr)^2
\right]
\\
&\le C_\star\rho_{n,d,\epsilon}.
\end{aligned}
\]
Together with the causal lower bound from the first step, these are the asserted causal inequalities. Since \(n,d,\epsilon\) were arbitrary in the stated range, all bounds hold with the chosen universal constants.
% lean: CausalSmith.Stat.OptvalueVanishingoverlapRate.matched_frontier
\end{enumerate}
\end{proof}

\begin{proof}[Proof of \cref{obj:thm:consistency-threshold}]
\begin{enumerate}
\item Fix the universal constants and estimator supplied by
\cref{obj:thm:matched-frontier}. For integers \(n\ge1\), \(d\ge2\), and
\(0<\epsilon\le1/2\), write
\[
L_d:=\log(ed),
\qquad
\rho_{n,d,\epsilon}
:=\min\left\{1,\frac{d}{n\epsilon L_d}\right\}.
\]
The supplied constants satisfy \(0<c_\star\le C_\star\), and the observed-risk bounds are
\[
c_\star\rho_{n,d,\epsilon}
\le \mathfrak R_{n,d,\epsilon}
\le
\sup_{\mathbb P\in\mathcal D_{d,\epsilon}^{\mathrm{obs}}}
E_{\mathbb P}\!\left[
\bigl(\widehat V_{n,d,\epsilon}^{\mathrm{AF}}-\Psi(\mathbb P)\bigr)^2
\right]
\le C_\star\rho_{n,d,\epsilon}.
\]
% lean: CausalSmith.Stat.OptvalueVanishingoverlapRate.matched_frontier

\item Apply the final equivalence in
\cref{obj:thm:weighted-separation-and-frontier} with the sample size equal to the sequence index. These sample sizes tend to infinity, and the stipulated alphabet and overlap conditions satisfy its remaining hypotheses. It gives precisely
\[
\begin{gathered}
\exists\,(\widehat V_n)_{n\in\mathbb N}
\text{ of measurable real-valued estimators based on \(n\) observations}:\\
\sup_{\mathbb P\in\mathcal D_{d_n,\epsilon_n}^{\mathrm{obs}}}
E_{\mathbb P}\!\left[
\bigl(\widehat V_n-\Psi(\mathbb P)\bigr)^2
\right]\longrightarrow0
\end{gathered}
\quad\Longleftrightarrow\quad
\frac{d_n}{n\epsilon_n L_{d_n}}\longrightarrow0.
\]
Here the ratio is evaluated for \(n\ge1\); the limit is unchanged by the initial index. At \(n=0\), the estimator may be assigned the constant value
\(\widehat V_0=0\). Since \(L_{d_n}=\log(ed_n)\), this proves the consistency assertion.
% lean: CausalSmith.Stat.OptvalueVanishingoverlapRate.weighted_separation_and_frontier

\item Fix an integer \(d\ge2\). The logarithmic bounds
\[
0\le\log d\le d-1
\]
follow respectively from \(d\ge1\) and the standard inequality
\(\log d\le d-1\) for \(d>0\). Consequently,
\[
L_d=1+\log d,
\qquad
0<L_d\le d,
\qquad
\lambda_d:=\frac{d}{L_d}\ge1.
\]
Choose
\[
c_d:=c_\star,
\qquad
C_d:=C_\star\frac{d}{L_d}.
\]
Because \(C_\star\ge c_\star>0\) and \(\lambda_d\ge1\),
\[
0<c_d=c_\star\le C_\star
\le C_\star\lambda_d=C_d.
\]

We will use the elementary comparison
\[
\min\{1,\xi_{\mathrm{cmp}}\}
\le \min\{1,\lambda_{\mathrm{cmp}}\xi_{\mathrm{cmp}}\}
\le \lambda_{\mathrm{cmp}}\min\{1,\xi_{\mathrm{cmp}}\},
\qquad
\lambda_{\mathrm{cmp}}\in[1,\infty),\quad
\xi_{\mathrm{cmp}}\in[0,\infty).
\]
For the first inequality,
\(\xi_{\mathrm{cmp}}\le\lambda_{\mathrm{cmp}}\xi_{\mathrm{cmp}}\), so monotonicity of the minimum gives
\[
\min\{1,\xi_{\mathrm{cmp}}\}
\le\min\{1,\lambda_{\mathrm{cmp}}\xi_{\mathrm{cmp}}\}.
\]
For the second inequality, if \(\xi_{\mathrm{cmp}}\le1\), then
\[
\min\{1,\lambda_{\mathrm{cmp}}\xi_{\mathrm{cmp}}\}
\le\lambda_{\mathrm{cmp}}\xi_{\mathrm{cmp}}
=\lambda_{\mathrm{cmp}}\min\{1,\xi_{\mathrm{cmp}}\}.
\]
If \(\xi_{\mathrm{cmp}}\ge1\), then
\[
\min\{1,\lambda_{\mathrm{cmp}}\xi_{\mathrm{cmp}}\}
\le1\le\lambda_{\mathrm{cmp}}
=\lambda_{\mathrm{cmp}}\min\{1,\xi_{\mathrm{cmp}}\}.
\]
These two cases include the common boundary \(\xi_{\mathrm{cmp}}=1\).

Now let \(n\ge1\) and \(0<\epsilon\le1/2\), and define
\[
\xi_{n,\epsilon}:=\frac{1}{n\epsilon}\ge0.
\]
The denominator is positive, and direct rearrangement gives
\[
\rho_{n,d,\epsilon}
=\min\{1,\lambda_d\xi_{n,\epsilon}\}.
\]
Applying the comparison with
\(\lambda_{\mathrm{cmp}}=\lambda_d\) and
\(\xi_{\mathrm{cmp}}=\xi_{n,\epsilon}\), the frontier bounds yield
\[
c_d\min\{1,\xi_{n,\epsilon}\}
\le c_\star\min\{1,\lambda_d\xi_{n,\epsilon}\}
\le\mathfrak R_{n,d,\epsilon}
\]
and
\[
\begin{aligned}
\mathfrak R_{n,d,\epsilon}
&\le
\sup_{\mathbb P\in\mathcal D_{d,\epsilon}^{\mathrm{obs}}}
E_{\mathbb P}\!\left[
\bigl(\widehat V_{n,d,\epsilon}^{\mathrm{AF}}-\Psi(\mathbb P)\bigr)^2
\right]\\
&\le C_\star\min\{1,\lambda_d\xi_{n,\epsilon}\}\\
&\le C_\star\lambda_d\min\{1,\xi_{n,\epsilon}\}
=C_d\min\left\{1,\frac{1}{n\epsilon}\right\}.
\end{aligned}
\]
Thus \(c=c_d\) and \(C=C_d\) give the fixed-alphabet bounds throughout the required range.
% lean: CausalSmith.Stat.OptvalueVanishingoverlapRate.logAlphabet_bounds
% lean: CausalSmith.Stat.OptvalueVanishingoverlapRate.min_one_mul_bounds
% lean: CausalSmith.Stat.OptvalueVanishingoverlapRate.consistency_threshold

\item Fix \(0<\epsilon\le1/2\), and choose
\[
c_\epsilon:=c_\star,
\qquad
C_\epsilon:=\frac{C_\star}{\epsilon}.
\]
Since \(\epsilon>0\) and \(\epsilon\le1/2\le1\),
\[
\lambda_\epsilon:=\frac1\epsilon\ge1.
\]
It follows that
\[
0<c_\epsilon=c_\star\le C_\star
\le C_\star\lambda_\epsilon=C_\epsilon.
\]
For integers \(n\ge1\) and \(d\ge2\), the logarithmic bounds established above give \(L_d>0\). Define
\[
\xi_{n,d}:=\frac{d}{nL_d}\ge0.
\]
Then
\[
\rho_{n,d,\epsilon}
=\min\{1,\lambda_\epsilon\xi_{n,d}\}.
\]
Applying the same elementary comparison with
\(\lambda_{\mathrm{cmp}}=\lambda_\epsilon\) and
\(\xi_{\mathrm{cmp}}=\xi_{n,d}\), we obtain
\[
c_\epsilon\min\{1,\xi_{n,d}\}
\le c_\star\min\{1,\lambda_\epsilon\xi_{n,d}\}
\le\mathfrak R_{n,d,\epsilon}
\]
and
\[
\begin{aligned}
\mathfrak R_{n,d,\epsilon}
&\le
\sup_{\mathbb P\in\mathcal D_{d,\epsilon}^{\mathrm{obs}}}
E_{\mathbb P}\!\left[
\bigl(\widehat V_{n,d,\epsilon}^{\mathrm{AF}}-\Psi(\mathbb P)\bigr)^2
\right]\\
&\le C_\star\min\{1,\lambda_\epsilon\xi_{n,d}\}\\
&\le C_\star\lambda_\epsilon\min\{1,\xi_{n,d}\}
=C_\epsilon\min\left\{1,\frac{d}{n\log(ed)}\right\}.
\end{aligned}
\]
Taking \(c=c_\epsilon\) and \(C=C_\epsilon\) proves the fixed-overlap bounds.
% lean: CausalSmith.Stat.OptvalueVanishingoverlapRate.consistency_threshold
\end{enumerate}
\end{proof}

\begin{proof}[Proof of \cref{obj:thm:observable-upper}]
\textit{1. Pilot calibration and auxiliary experiment.}
Choose
\[
H_0=1025.
\]
Thus \(H_0>0\), \(H_0\ge1\), and \(H_0\ge1024\), the universal self-normalization multiplier.

We first establish the auxiliary risk bound for integers \(n,d\ge1\), an overlap parameter \(0<\epsilon\le1/2\), and a law \(\mathbb P\in\mathcal D_{d,\epsilon}^{\mathrm{obs}}\). Write
\[
L=L_d=\log(ed),\qquad m=\frac n8,\qquad
\delta_{\mathrm p}=\frac Lm,\qquad
q_x=(q_{ay,x})_{a,y\in\{0,1\}}.
\]
The atom probabilities are those of \cref{obj:def:observed-class}; in particular, \(m>0\) and \(L=1+\log d\ge1\).

Consider the auxiliary experiment that draws \(\mathsf M\sim\operatorname{Pois}(n/4)\) and, conditional on \(\mathsf M\), draws \(\mathsf M\) independent observations with law \(\mathbb P\), together with independent fair marks. Its law on the disjoint union of all finite marked samples is
\[
\mathcal Q_{\mathbb P}
=
\sum_{k=0}^{\infty}
e^{-n/4}\frac{(n/4)^k}{k!}
\left(\mathbb P\otimes\operatorname{Bernoulli}(1/2)\right)^{\otimes k}.
\]
Write \(\mathbb E_{\mathrm{aux}}\) for expectation under this law. Construct the counts and cell statistics of \cref{obj:def:armwise-estimator} for every finite marked sample, with the auxiliary count unrestricted. For arrays
\(\boldsymbol t^{\mathrm p},\boldsymbol t^{\mathrm e}\in[0,1]^{\mathcal J_d}\), indexed by the observed atom set defined in \cref{obj:synth_1}, the count generating function is
\[
\mathbb E_{\mathrm{aux}}
\prod_{x,a,y}
(t^{\mathrm p}_{ay,x})^{N'_{ay,x}}
(t^{\mathrm e}_{ay,x})^{N_{ay,x}}
=
\exp\left\{
m\sum_{x,a,y}q_{ay,x}(t^{\mathrm p}_{ay,x}-1)
+
m\sum_{x,a,y}q_{ay,x}(t^{\mathrm e}_{ay,x}-1)
\right\}.
\]
Indeed, conditional on the auxiliary count, each marked observation contributes one category factor; summing over the Poisson count gives the displayed exponential. Consequently,
\[
N'_{ay,x},N_{ay,x}\sim\operatorname{Pois}(mq_{ay,x}),
\]
and all these counts are mutually independent. In particular, the complete cell statistics \(T_x\) are independent across distinct cells.
% lean: CausalSmith.Stat.OptvalueVanishingoverlapRate.markedCount_poisson_hasLaw
% lean: CausalSmith.Stat.OptvalueVanishingoverlapRate.markedPoissonCellValue_pairwise_indepFun

\textit{2. Polynomial oscillation and coefficient control.}
Use the centers, radii, rectangles, polynomial, and clipping scales of \cref{obj:def:armwise-estimator}. Every coordinate radius is positive, since
\[
h_{ay,x}=H_0\left(\sqrt{c_{ay,x}\delta_{\mathrm p}}+\delta_{\mathrm p}\right)>0,
\qquad
\frac{h_{ay,x}}2\le r_{ay,x}\le h_{ay,x}.
\]
Moreover,
\[
c_{ay,x}\le b_{ay,x}\le c_{ay,x}+h_{ay,x}.
\]
Thus \(b_x\ne0\). Applying \cref{obj:lem:armwise-modulus} with comparison vector \(b_x\), for every \(v\in Q_x\), gives
\[
|F_\epsilon(v)-F_\epsilon(b_x)|
\le
2\sum_{a=0}^1
\left(1+\frac{s(b_x)}{D_a(b_x)}\right)R_{a,x}
=S_x^{\mathrm{aw}}.
\]
The Jackson convolution is positive and normalized, so the same bound holds for its polynomial:
\[
\sup_{v\in Q_x}
|P_{x,K_n}(v)-F_\epsilon(b_x)|
\le S_x^{\mathrm{aw}}.
\]

For completeness, define the Chebyshev polynomials by
\[
\mathsf{Ch}_0(\xi_\circ)=1,\qquad
\mathsf{Ch}_1(\xi_\circ)=\xi_\circ,\qquad
\mathsf{Ch}_{j+1}(\xi_\circ)
=2\xi_\circ\mathsf{Ch}_j(\xi_\circ)-\mathsf{Ch}_{j-1}(\xi_\circ),
\quad j\ge1.
\]
Here \(\xi_\circ\) is a real scalar. The recurrence implies
\[
\sum_{\ell=0}^{j}
\left|[\xi_\circ^\ell]\mathsf{Ch}_j(\xi_\circ)\right|
\le(1+\sqrt2)^j,
\]
by induction and the identity
\((1+\sqrt2)^2=2(1+\sqrt2)+1\).

With \(D=2(K_n-1)\), expand the normalized polynomial as
\[
P_{x,K_n}(b_x+r_x\odot\xi)-F_\epsilon(b_x)
=
\sum_{\boldsymbol\nu\in\{0,\ldots,D\}^4}
\gamma_{x,\boldsymbol\nu}
\prod_{a,y}\mathsf{Ch}_{\nu_{ay}}(\xi_{ay}),
\qquad \xi\in[-1,1]^4.
\]
Cosine orthogonality expresses each coefficient as an integral of the bounded polynomial against cosine factors, with normalization factor at most \(16\). Hence
\[
|\gamma_{x,\boldsymbol\nu}|\le16S_x^{\mathrm{aw}}.
\]
Converting the tensor Chebyshev expansion to the monomial expansion defining \(\beta_{x,\boldsymbol\alpha}\) yields
\[
\sum_{\boldsymbol\alpha}|\beta_{x,\boldsymbol\alpha}|
\le
16(D+1)^4(1+\sqrt2)^{4D}S_x^{\mathrm{aw}}.
\]
The numerical factor satisfies
\[
\left\{16(D+1)^4(1+\sqrt2)^{4D}\right\}^2
\le e^{256+24D}\le e^{304K_n}.
\]
Indeed, \(D+1\le e^D\), \(1+\sqrt2\le e^2\), \(256\le e^{256}\), and \(D\le2K_n\), \(K_n\ge1\).
% lean: CausalSmith.Stat.OptvalueVanishingoverlapRate.jacksonCoeff_sum_abs_le
% lean: CausalSmith.Stat.OptvalueVanishingoverlapRate.jacksonCoeff_factor_square_le_exp

\textit{3. Factorial second moments and universal degree tuning.}
Define the successful-localization event by
\[
\mathcal G_x=\{q_x\in Q_x\}.
\]
On this event,
\[
|q_{ay,x}-b_{ay,x}|\le r_{ay,x},
\qquad
\frac{q_{ay,x}}{m r_{ay,x}^2}\le\frac8L.
\]
To verify the second inequality, the coordinate formulas give
\[
c_{ay,x}\delta_{\mathrm p}\le h_{ay,x}^2,\qquad
\delta_{\mathrm p}\le h_{ay,x},\qquad
r_{ay,x}\ge\frac{h_{ay,x}}2,
\]
because \(H_0\ge1\). Since \(q_{ay,x}\le c_{ay,x}+h_{ay,x}\),
\[
q_{ay,x}\delta_{\mathrm p}
\le(c_{ay,x}+h_{ay,x})\delta_{\mathrm p}
\le2h_{ay,x}^2
\le8r_{ay,x}^2,
\]
which proves the claim.

For a Poisson count \(\mathsf N_\circ\) with mean \(m q_\circ\), where \(q_\circ\ge0\), the factorial identities are
\[
\mathbb E(\mathsf N_\circ)_j=(m q_\circ)^j,
\]
\[
(\mathsf N_\circ)_j(\mathsf N_\circ)_k
=
\sum_{\ell=0}^{\min(j,k)}
\binom j\ell\binom k\ell\ell!\,
(\mathsf N_\circ)_{j+k-\ell},
\qquad j,k\in\mathbb N.
\]
The first follows by shifting the Poisson series. The second counts the overlap of two ordered selections of distinct indices. Substituting the first identity into the centered factorial statistic of \cref{obj:def:armwise-estimator}, for every real \(\zeta\), and applying the binomial theorem gives
\[
\mathbb E U_j(\mathsf N_\circ;\zeta)
=
\sum_{j_{\mathrm{raw}}=0}^{j}
\binom j{j_{\mathrm{raw}}}(-\zeta)^{j-j_{\mathrm{raw}}}
q_\circ^{j_{\mathrm{raw}}}
=(q_\circ-\zeta)^j.
\]
For the mixed second moment, substitution of the raw factorial product identity gives the finite triple sum
\[
\begin{aligned}
\mathbb E[U_j(\mathsf N_\circ;\zeta)U_k(\mathsf N_\circ;\zeta)]
={}&
\sum_{j_{\mathrm{raw}}=0}^{j}
\sum_{k_{\mathrm{raw}}=0}^{k}
\binom j{j_{\mathrm{raw}}}\binom k{k_{\mathrm{raw}}}
(-\zeta)^{j-j_{\mathrm{raw}}+k-k_{\mathrm{raw}}}\\
&\quad\times
\sum_{\ell=0}^{\min(j_{\mathrm{raw}},k_{\mathrm{raw}})}
\binom{j_{\mathrm{raw}}}\ell
\binom{k_{\mathrm{raw}}}\ell\ell!
\left(\frac{q_\circ}{m}\right)^\ell
q_\circ^{j_{\mathrm{raw}}+k_{\mathrm{raw}}-2\ell}.
\end{aligned}
\]
Fixing the overlap order and regrouping the finite sums factors this expression as
\[
\begin{aligned}
\mathbb E[U_j(\mathsf N_\circ;\zeta)U_k(\mathsf N_\circ;\zeta)]
={}&\sum_{\ell=0}^{\min(j,k)}
\ell!\left(\frac{q_\circ}{m}\right)^\ell\\
&\times\left\{
\sum_{j_{\mathrm{raw}}=\ell}^{j}
\binom j{j_{\mathrm{raw}}}\binom{j_{\mathrm{raw}}}\ell
(-\zeta)^{j-j_{\mathrm{raw}}}
q_\circ^{j_{\mathrm{raw}}-\ell}
\right\}\\
&\times\left\{
\sum_{k_{\mathrm{raw}}=\ell}^{k}
\binom k{k_{\mathrm{raw}}}\binom{k_{\mathrm{raw}}}\ell
(-\zeta)^{k-k_{\mathrm{raw}}}
q_\circ^{k_{\mathrm{raw}}-\ell}
\right\}.
\end{aligned}
\]
Terms with a raw order smaller than \(\ell\) vanish because the corresponding binomial coefficient is zero. For the remaining terms, the shifted binomial identity is
\[
\begin{aligned}
\sum_{j_{\mathrm{raw}}=\ell}^{j}
\binom j{j_{\mathrm{raw}}}\binom{j_{\mathrm{raw}}}\ell
(-\zeta)^{j-j_{\mathrm{raw}}}
q_\circ^{j_{\mathrm{raw}}-\ell}
&=
\binom j\ell
\sum_{i_{\mathrm{shift}}=0}^{j-\ell}
\binom{j-\ell}{i_{\mathrm{shift}}}
(-\zeta)^{j-\ell-i_{\mathrm{shift}}}
q_\circ^{i_{\mathrm{shift}}}\\
&=\binom j\ell(q_\circ-\zeta)^{j-\ell},
\qquad 0\le\ell\le j.
\end{aligned}
\]
The first equality substitutes \(j_{\mathrm{raw}}=\ell+i_{\mathrm{shift}}\) and uses
\[
\binom j{\ell+i_{\mathrm{shift}}}
\binom{\ell+i_{\mathrm{shift}}}\ell
=
\binom j\ell\binom{j-\ell}{i_{\mathrm{shift}}},
\qquad
 i_{\mathrm{shift}}\in\{0,\ldots,j-\ell\};
\]
the second is the binomial theorem. Applying this identity to both factors yields
\[
\mathbb E[U_j(\mathsf N_\circ;\zeta)U_k(\mathsf N_\circ;\zeta)]
=
\sum_{\ell=0}^{\min(j,k)}
\binom j\ell\binom k\ell\ell!
\left(\frac{q_\circ}{m}\right)^\ell
(q_\circ-\zeta)^{j+k-2\ell}.
\]
Taking \(k=j\) therefore gives
\[
\mathbb E U_j(\mathsf N_\circ;\zeta)^2
=
\sum_{\ell=0}^{j}
\binom j\ell^2\ell!
\left(\frac{q_\circ}{m}\right)^\ell
(q_\circ-\zeta)^{2j-2\ell}.
\]
Thus, conditional on a pilot in \(\mathcal G_x\), for \(0\le j\le D\),
\[
\mathbb E_{\mathrm{aux}}
\left[
\left.
\frac{U_j(N_{ay,x};b_{ay,x})^2}{r_{ay,x}^{2j}}
\right|N'
\right]
\le
\sum_{\ell=0}^{j}\frac{j^{2\ell}}{\ell!}
\left(\frac8L\right)^\ell
\le e^{8j^2/L}.
\]
Here \(N'\) denotes the complete pilot-count array. Evaluation-coordinate independence therefore gives
\[
\mathbb E_{\mathrm{aux}}[Z_x\mid N']
=P_{x,K_n}(q_x),
\]
\[
\mathbb E_{\mathrm{aux}}
[(Z_x-F_\epsilon(b_x))^2\mid N']
\le
\left(\sum_{\boldsymbol\alpha}|\beta_{x,\boldsymbol\alpha}|\right)^2
e^{32D^2/L}
\quad\text{on }\mathcal G_x.
\]
For the last inequality, each normalized factorial product has second moment at most \(e^{32D^2/L}\); Cauchy--Schwarz bounds each cross term in the polynomial square by this same envelope.

Choose the numerical constants
\[
\kappa=\frac1{17952},
\qquad
C_{\mathrm{fac}}=\exp\left(1632+\frac1{16}\right).
\]
This choice satisfies
\[
304\kappa+256\kappa^2\le\frac1{16}.
\]
The degree definition gives, for every \(d\ge1\),
\[
\frac{\kappa L}{2}\le K_n\le2+\kappa L.
\]
For the lower bound, split according as \(\kappa L\le4\) or \(\kappa L>4\), and use respectively \(K_n\ge2\) or
\(\lfloor\kappa L\rfloor>\kappa L-1\). Also,
\[
K_n^2\le8+2\kappa^2L^2.
\]
Combining these inequalities with the coefficient estimate gives
\[
\begin{aligned}
\left(\sum_{\boldsymbol\alpha}|\beta_{x,\boldsymbol\alpha}|\right)^2
e^{32D^2/L}
&\le
(S_x^{\mathrm{aw}})^2
e^{304K_n+128K_n^2/L}\\
&\le
(S_x^{\mathrm{aw}})^2e^{1632+L/16}\\
&=
C_{\mathrm{fac}}d^{1/16}(S_x^{\mathrm{aw}})^2.
\end{aligned}
\]
Thus the degree tuning is uniform over the entire auxiliary alphabet range.
% lean: CausalSmith.Stat.OptvalueVanishingoverlapRate.pilotFactorialCellValue_square_envelope
% lean: CausalSmith.Stat.OptvalueVanishingoverlapRate.pilotJacksonCoeff_square_tuning_exists

\textit{4. Clipping on successful pilots.}
On \(\mathcal G_x\), the polynomial oscillation bound places \(P_{x,K_n}(q_x)\) in the clipping interval, since \(d^{1/4}\ge1\). Projection onto that interval consequently gives
\[
\mathbb E_{\mathrm{aux}}
[(T_x-P_{x,K_n}(q_x))^2\mid N']
\le
C_{\mathrm{fac}}d^{1/16}(S_x^{\mathrm{aw}})^2.
\]
Indeed, projection contracts distance from \(P_{x,K_n}(q_x)\), and
\[
\mathbb E_{\mathrm{aux}}
[(Z_x-P_{x,K_n}(q_x))^2\mid N']
\le
\mathbb E_{\mathrm{aux}}
[(Z_x-F_\epsilon(b_x))^2\mid N'].
\]
For the clipping bias, the scalar projection inequality
\[
|\Pi_{[-\mathfrak u_{\mathrm{clip}},\mathfrak u_{\mathrm{clip}}]}(\mathfrak v_{\mathrm{clip}})-\mathfrak v_{\mathrm{clip}}|
\le\frac{\mathfrak v_{\mathrm{clip}}^2}{\mathfrak u_{\mathrm{clip}}},
\qquad \mathfrak u_{\mathrm{clip}}>0,\quad
\mathfrak v_{\mathrm{clip}}\in\mathbb R,
\]
follows by considering \(|\mathfrak v_{\mathrm{clip}}|\le\mathfrak u_{\mathrm{clip}}\) and \(|\mathfrak v_{\mathrm{clip}}|>\mathfrak u_{\mathrm{clip}}\). Applying it with clipping radius \(d^{1/4}S_x^{\mathrm{aw}}>0\), and using factorial unbiasedness, yields
\[
\left|
\mathbb E_{\mathrm{aux}}[T_x\mid N']-P_{x,K_n}(q_x)
\right|
\le C_{\mathrm{fac}}d^{-3/16}S_x^{\mathrm{aw}}
\quad\text{on }\mathcal G_x.
\]
% lean: CausalSmith.Stat.OptvalueVanishingoverlapRate.pilotClippedCellValue_tuning_exists

\textit{5. Approximation error in occupied cells.}
For \(p_x>0\), define
\[
\rho_{1x}=\pi_x,\qquad \rho_{0x}=1-\pi_x.
\]
By \cref{obj:def:observed-class,obj:ass:shrinking-overlap},
\[
s(q_x)=p_x,\qquad
s_a(q_x)=p_x\rho_{ax}>0,\qquad
D_a(q_x)=s_a(q_x),\qquad
\rho_{ax}\ge\epsilon.
\]
It follows from \cref{obj:def:armwise-extension,obj:def:observed-value} that
\[
F_\epsilon(q_x)=p_x\max_a\mu_{ax},
\qquad
\Psi(\mathbb P)=\sum_{x\in[d]}F_\epsilon(q_x).
\]
The same cell identity holds when \(p_x=0\), because then \(q_x=0\).

We record the kernel moments used in the approximation. Define, for \(t\in[-\pi,\pi]\), the continuously extended sine ratio and its fourth-power mass by
\[
\mathcal S_{K_n}(t)=
\begin{cases}
\displaystyle\frac{\sin(K_nt/2)}{\sin(t/2)},&t\ne0,\\[4pt]
K_n,&t=0,
\end{cases}
\qquad
\mathcal A_{K_n}=\int_{-\pi}^{\pi}\mathcal S_{K_n}(t)^4\,dt.
\]
The finite cosine expansion is
\[
\mathcal S_{K_n}(t)^2
=K_n+2\sum_{j=1}^{K_n-1}(K_n-j)\cos(jt).
\]
Indeed, multiplying the right side by \(2(1-\cos t)\) and using the cosine product identity telescopes to \(2(1-\cos(K_nt))\); division by \(4\sin^2(t/2)\) proves the identity for \(t\ne0\), and continuity supplies its value at zero. Cosine orthogonality and the finite square sum give
\[
\sum_{j=1}^{K_n-1}\{2(K_n-j)\}^2
=4\sum_{j=1}^{K_n-1}j^2
=\frac23(K_n-1)K_n(2K_n-1),
\]
where the square-sum formula follows by summing the successive differences of cubes. Consequently,
\[
\begin{aligned}
\mathcal A_{K_n}
&=2\pi K_n^2+
\pi\sum_{j=1}^{K_n-1}\{2(K_n-j)\}^2\\
&=\frac{2\pi}{3}K_n(2K_n^2+1)
\ge\frac{4\pi}{3}K_n^3,
\qquad
J_{K_n}(t)=\frac{\mathcal S_{K_n}(t)^4}{\mathcal A_{K_n}}.
\end{aligned}
\]
The same cosine expansion, integrated without squaring, yields
\[
\int_{-\pi}^{\pi}\mathcal S_{K_n}(t)^2\,dt=2\pi K_n.
\]
For \(|t|\le\pi\), the sine lower bound gives
\[
|t|\le\pi|\sin(t/2)|,
\qquad
t^2\le\pi^2\sin^2(t/2).
\]
For \(t\ne0\), it follows that
\[
\begin{aligned}
t^2\mathcal S_{K_n}(t)^4
&\le\pi^2\sin^2(K_nt/2)\mathcal S_{K_n}(t)^2\\
&\le\pi^2\mathcal S_{K_n}(t)^2.
\end{aligned}
\]
At \(t=0\), this inequality follows directly from its zero left side. Integrating and normalizing therefore gives
\[
\int_{-\pi}^{\pi}t^2\mathcal S_{K_n}(t)^4\,dt
\le2\pi^3K_n,
\qquad
\int_{-\pi}^{\pi}t^2J_{K_n}(t)\,dt
\le\frac{24}{K_n^2}\le\frac{64}{K_n^2},
\]
using the mass lower bound and \(\pi^2\le16\).
Finally, the nonnegative square \((K_n|t|-2)^2\) implies
\[
|t|\le\frac{K_n}{4}t^2+\frac1{K_n}.
\]
Since the kernel is nonnegative and integrates to one, we obtain
\[
\begin{aligned}
\int_{-\pi}^{\pi}|t|J_{K_n}(t)\,dt
&\le\frac{K_n}{4}
\int_{-\pi}^{\pi}t^2J_{K_n}(t)\,dt+\frac1{K_n}\\
&\le\frac7{K_n}\le\frac{32}{K_n}.
\end{aligned}
\]
% lean: Causalean.Mathlib.Analysis.JacksonApproximation.jrawMass_eq
% lean: Causalean.Mathlib.Analysis.JacksonApproximation.jackson_second_moment
% lean: Causalean.Mathlib.Analysis.JacksonApproximation.jackson_first_moment
For \(q_{ay,x}=b_{ay,x}+r_{ay,x}\cos\theta_{ay,x}\), define
\[
\theta_{ay,x}
=
\arccos\left(\frac{q_{ay,x}-b_{ay,x}}{r_{ay,x}}\right)
\quad\text{on }\mathcal G_x.
\]
Taylor's formula gives
\[
\begin{aligned}
r_{ay,x}
|\cos(\theta_{ay,x}+t)-\cos\theta_{ay,x}|
\le{}&
\sqrt{(q_{ay,x}-\ell_{ay,x})(u_{ay,x}-q_{ay,x})}\,|t|\\
&+\frac{r_{ay,x}}2t^2.
\end{aligned}
\]
Applying \cref{obj:lem:armwise-modulus} at \(q_x\) inside the tensor convolution therefore gives
\[
\begin{aligned}
|P_{x,K_n}(q_x)-F_\epsilon(q_x)|
\le64\sum_{a=0}^1
\left(1+\frac{p_x}{s_a(q_x)}\right)
\left\{
\frac{\sum_y\sqrt{(q_{ay,x}-\ell_{ay,x})(u_{ay,x}-q_{ay,x})}}{K_n}
+\frac{R_{a,x}}{K_n^2}
\right\}.
\end{aligned}
\]

The localized coordinate formulas imply
\[
h_{ay,x}
\le H_0(H_0+2)
\left(\sqrt{q_{ay,x}\delta_{\mathrm p}}+\delta_{\mathrm p}\right).
\]
To see this, the lower-endpoint inequality implies
\[
c_{ay,x}-H_0\left(\sqrt{c_{ay,x}\delta_{\mathrm p}}+\delta_{\mathrm p}\right)
\le q_{ay,x}.
\]
Solving the resulting quadratic inequality in \(\sqrt{c_{ay,x}}\) gives
\[
\sqrt{c_{ay,x}\delta_{\mathrm p}}
\le
\sqrt{q_{ay,x}\delta_{\mathrm p}}+(H_0+1)\delta_{\mathrm p},
\]
which yields the stated half-width bound.

There is also the boundary-sensitive estimate
\[
\sqrt{(q_{ay,x}-\ell_{ay,x})(u_{ay,x}-q_{ay,x})}
\le2H_0(H_0+2)\sqrt{q_{ay,x}\delta_{\mathrm p}}.
\]
For \(q_{ay,x}\le\delta_{\mathrm p}\), bound the left distance by \(q_{ay,x}\), the right distance by \(2h_{ay,x}\), and the half width by
\(2H_0(H_0+2)\delta_{\mathrm p}\). For \(q_{ay,x}>\delta_{\mathrm p}\), the product of the two distances is at most \(h_{ay,x}^2\), and the half-width bound applies. Summing coordinate radii also gives
\[
R_{a,x}
\le2H_0(H_0+2)
\left(\sqrt{s_a(q_x)\delta_{\mathrm p}}+\delta_{\mathrm p}\right).
\]

For \(d\ge2\), the mass bound in \cref{obj:lem:armwise-modulus} supplies
\[
\sum_{a=0}^1\frac{p_x}{s_a(q_x)}
\sum_{y=0}^1\sqrt{q_{ay,x}}
\le2\sqrt2\sqrt{\frac{p_x}{\epsilon}}.
\]
For \(d=1\), the unique cell has \(p_x=1\), and the same bound follows directly from the two arm masses. For each arm, Cauchy--Schwarz and \(s_a(q_x)\ge\epsilon p_x>0\) give
\[
\sum_{y=0}^1\sqrt{q_{ay,x}}\le\sqrt{2s_a(q_x)},
\qquad
\frac{p_x}{s_a(q_x)}\sum_{y=0}^1\sqrt{q_{ay,x}}
\le\frac{\sqrt2\,p_x}{\sqrt{s_a(q_x)}}
\le\sqrt{\frac{2p_x}{\epsilon}}.
\]
Adding the two arm inequalities proves the displayed mass bound also in this case.
% lean: CausalSmith.Stat.OptvalueVanishingoverlapRate.observed_armwise_sqrt_bound
Together with \(1+p_x/s_a(q_x)\le2p_x/s_a(q_x)\), this bound now yields
\[
\begin{aligned}
|P_{x,K_n}(q_x)-F_\epsilon(q_x)|
\le512H_0(H_0+2)
\left\{
\frac{\sqrt2}{K_n}\sqrt{\frac{p_xL}{m\epsilon}}
+
\frac1{K_n^2}
\left(\sqrt{\frac{p_xL}{m\epsilon}}+\frac L{m\epsilon}\right)
\right\}.
\end{aligned}
\]
Define
\[
\eta_x=
\sqrt{\frac{p_x}{m\epsilon L}}+\frac1{m\epsilon L},
\]
\[
C_{\mathrm{occ}}
=
512H_0(H_0+2)\sqrt2
\left(\frac1{\kappa/2}+\frac1{(\kappa/2)^2}\right).
\]
Using \(K_n\ge(\kappa/2)L\) and \(L\ge1\), we conclude that
\[
|P_{x,K_n}(q_x)-F_\epsilon(q_x)|
\le C_{\mathrm{occ}}\eta_x
\quad\text{on }\mathcal G_x,\quad p_x>0.
\]
Combining this estimate with the clipping bounds and integrating over the pilot gives
\[
\mathbb E_{\mathrm{aux}}
[(T_x-F_\epsilon(q_x))^2\mathbf1_{\mathcal G_x}]
\le
2C_{\mathrm{fac}}d^{1/16}
\mathbb E_{\mathrm{aux}}(S_x^{\mathrm{aw}})^2
+2C_{\mathrm{occ}}^2\eta_x^2,
\]
\[
\left|
\mathbb E_{\mathrm{aux}}
[(T_x-F_\epsilon(q_x))\mathbf1_{\mathcal G_x}]
\right|
\le
C_{\mathrm{occ}}\eta_x
+
C_{\mathrm{fac}}d^{-3/16}
\mathbb E_{\mathrm{aux}}S_x^{\mathrm{aw}}.
\]
The squared bound uses
\((v+w)^2\le2v^2+2w^2\); the mean bound uses the triangle inequality.
% lean: CausalSmith.Stat.OptvalueVanishingoverlapRate.observed_pilotJackson_mean_bias_normalized_le
% lean: CausalSmith.Stat.OptvalueVanishingoverlapRate.markedPoissonCellValue_occupied_good_moments_tuning_exists

\textit{6. Null cells and uniform successful-pilot bounds.}
If \(p_x=0\), every pilot and evaluation count in that cell is zero almost surely. Hence
\[
q_x=0,\qquad
b_{ay,x}=r_{ay,x}=\frac{H_0L}{2m},\qquad
\ell_{ay,x}=0,
\qquad
\mathbb P_{\mathrm{aux}}(\mathcal G_x)=1.
\]
At the lower corner, the factorial statistic equals the polynomial value:
\[
Z_x=P_{x,K_n}(0).
\]
This follows either by substituting zero counts in the factorial expansion or by its expectation identity, since the evaluation counts are constant. The polynomial oscillation bound places this value inside the clipping interval, so
\[
T_x=P_{x,K_n}(0).
\]
From \cref{obj:def:armwise-extension},
\[
0\le F_\epsilon(v)\le s(v),\qquad v\in\mathbb R_+^4.
\]
At the null corner, a displaced coordinate in the convolution is
\(r_{ay,x}(1-\cos t)\le r_{ay,x}t^2/2\). The second kernel moment thus gives
\[
|T_x-F_\epsilon(0)|
\le\frac{32}{K_n^2}\sum_{a,y}r_{ay,x}
=\frac{64H_0L}{mK_n^2}.
\]
Define
\[
C_{\mathrm{null}}=\frac{64H_0}{(\kappa/2)^2},
\qquad
C_{\mathrm{app}}=C_{\mathrm{occ}}+C_{\mathrm{null}}.
\]
Since \(K_n\ge(\kappa/2)L\) and \(\epsilon\le1\), the null-cell error is at most \(C_{\mathrm{null}}\eta_x\).

Therefore, for every cell, occupied or null, the successful-pilot estimates take the common form
\[
\mathbb E_{\mathrm{aux}}
[(T_x-F_\epsilon(q_x))^2\mathbf1_{\mathcal G_x}]
\le
C_{\mathrm v}d^{1/16}
\mathbb E_{\mathrm{aux}}(S_x^{\mathrm{aw}})^2
+C_{\mathrm e}\eta_x^2,
\]
\[
\left|
\mathbb E_{\mathrm{aux}}
[(T_x-F_\epsilon(q_x))\mathbf1_{\mathcal G_x}]
\right|
\le
C_{\mathrm b}
\left(\eta_x+d^{-3/16}\mathbb E_{\mathrm{aux}}S_x^{\mathrm{aw}}\right),
\]
where the universal constants are
\[
C_{\mathrm v}=2C_{\mathrm{fac}},\qquad
C_{\mathrm e}=2C_{\mathrm{app}}^2,\qquad
C_{\mathrm b}=C_{\mathrm{app}}+C_{\mathrm{fac}}.
\]
These enlargements use the nonnegativity of the clipping scales and of \(\eta_x\).
% lean: CausalSmith.Stat.OptvalueVanishingoverlapRate.markedPoissonCellValue_null_good_moments_le
% lean: CausalSmith.Stat.OptvalueVanishingoverlapRate.markedPoissonStatistic_active_rate_tuning_exists

\textit{7. Unconditional clipping-scale moments.}
For each pilot, the coordinate inequalities \(r_{ay,x}\le h_{ay,x}\) and \(c_{ay,x}\le b_{ay,x}\) imply
\[
R_{a,x}
\le2H_0\left(\sqrt{s_a(b_x)\delta_{\mathrm p}}+\delta_{\mathrm p}\right).
\]
Using \(D_a(b_x)\ge s_a(b_x)\) and
\(D_a(b_x)\ge\epsilon s(b_x)\), separately according as
\(s_a(b_x)\ge\epsilon s(b_x)\) or
\(s_a(b_x)<\epsilon s(b_x)\), gives
\[
\frac{s(b_x)}{D_a(b_x)}R_{a,x}
\le
2H_0\left(
\sqrt{\frac{s(b_x)\delta_{\mathrm p}}{\epsilon}}
+\frac{\delta_{\mathrm p}}{\epsilon}
\right).
\]
The unweighted radius satisfies this same bound, since
\(s_a(b_x)\le s(b_x)\) and \(\epsilon\le1\). Consequently,
\[
S_x^{\mathrm{aw}}
\le16H_0\left(
\sqrt{\frac{s(b_x)\delta_{\mathrm p}}{\epsilon}}
+\frac{\delta_{\mathrm p}}{\epsilon}
\right),
\]
\[
(S_x^{\mathrm{aw}})^2
\le512H_0^2\left(
\frac{s(b_x)\delta_{\mathrm p}}{\epsilon}
+\frac{\delta_{\mathrm p}^2}{\epsilon^2}
\right).
\]

The midpoint formulas also give
\[
b_{ay,x}\le(1+H_0)c_{ay,x}+2H_0\delta_{\mathrm p},
\qquad
s(b_x)\le(1+H_0)\sum_{a,y}c_{ay,x}+8H_0\delta_{\mathrm p}.
\]
For example, use \(b_{ay,x}\le c_{ay,x}+h_{ay,x}\) and
\(\sqrt{c_{ay,x}\delta_{\mathrm p}}\le c_{ay,x}+\delta_{\mathrm p}\).
Because \(\mathbb E_{\mathrm{aux}}c_{ay,x}=q_{ay,x}\), integration yields
\[
\mathbb E_{\mathrm{aux}}(S_x^{\mathrm{aw}})^2
\le512H_0^2
\left\{
(1+H_0)\frac{p_x\delta_{\mathrm p}}{\epsilon}
+8H_0\frac{\delta_{\mathrm p}^2}{\epsilon}
+\frac{\delta_{\mathrm p}^2}{\epsilon^2}
\right\}.
\]
Define, for every cell,
\[
\mathfrak a_x
=
\frac{p_xL}{m\epsilon}
+\frac{L^2}{m^2\epsilon^2},
\qquad
C_{\mathrm s}=512H_0^2(1+9H_0).
\]
Since \(\epsilon\le1\),
\[
\mathbb E_{\mathrm{aux}}(S_x^{\mathrm{aw}})^2
\le C_{\mathrm s}\mathfrak a_x.
\]
% lean: CausalSmith.Stat.OptvalueVanishingoverlapRate.markedPoissonClippingScale_integral_sq_le

\textit{8. Failed-pilot moments.}
Define the nonnegative coordinate scores by
\[
e_{ay,x}^{\mathrm p}
=
|c_{ay,x}-q_{ay,x}|+h_{ay,x},
\qquad
e_x^{\mathrm p}=\sum_{a,y}e_{ay,x}^{\mathrm p}.
\]
We need both a probability bound and moments retaining the cell mass:
\[
\mathbb P_{\mathrm{aux}}(\mathcal G_x^c)\le8e^{-40L},
\]
\[
\mathbb E_{\mathrm{aux}}(e_{ay,x}^{\mathrm p})^4
\le
\left(\frac{H_0}{1024}\right)^4
10^{20}
\left(\sqrt{q_{ay,x}\delta_{\mathrm p}}+\delta_{\mathrm p}\right)^4,
\]
\[
\mathbb E_{\mathrm{aux}}
[e_x^{\mathrm p}\mathbf1_{\mathcal G_x^c}]
\le
\frac{H_0}{1024}\,10^{160}e^{-20L}
\left(\sqrt{p_x\delta_{\mathrm p}}+\delta_{\mathrm p}\right).
\]

Here are the underlying calculations. For any
\(\lambda_\circ\ge0\) and
\(\mathsf N_\circ\sim\operatorname{Pois}(\lambda_\circ)\), define
\[
\mathfrak s_\circ=\sqrt{\lambda_\circ L}+L,\qquad
\mathfrak e_\circ
=
|\mathsf N_\circ-\lambda_\circ|
+1024\left(\sqrt{\mathsf N_\circ L}+L\right),
\qquad
\vartheta_\circ=\frac1{2\mathfrak s_\circ}.
\]
The centered Poisson generating function is
\[
\mathbb E e^{t(\mathsf N_\circ-\lambda_\circ)}
=
\exp\{\lambda_\circ(e^t-1-t)\}.
\]
If \(\lambda_\circ=0\), then \(\mathsf N_\circ=0\) almost surely, and both tail events below have probability zero. For \(\lambda_\circ>0\), define the deviation threshold and exponential tilt at an arbitrary nonnegative level by
\[
z_{\mathrm{tail}}\in[0,\infty),\qquad
\mathfrak x_{\mathrm{up}}
=\sqrt{2\lambda_\circ z_{\mathrm{tail}}}+2z_{\mathrm{tail}},
\qquad
\vartheta_{\mathrm{up}}
=\log\left(1+\frac{\mathfrak x_{\mathrm{up}}}{\lambda_\circ}\right)\ge0.
\]
Exponential Markov inequality and the generating function give
\[
\begin{aligned}
\Pr\{\mathsf N_\circ-\lambda_\circ>\mathfrak x_{\mathrm{up}}\}
&\le
\exp\left\{
-\vartheta_{\mathrm{up}}\mathfrak x_{\mathrm{up}}
+\lambda_\circ(e^{\vartheta_{\mathrm{up}}}-1-\vartheta_{\mathrm{up}})
\right\}\\
&=
\exp\left\{
-\bigl[(\lambda_\circ+\mathfrak x_{\mathrm{up}})
\vartheta_{\mathrm{up}}-\mathfrak x_{\mathrm{up}}\bigr]
\right\}.
\end{aligned}
\]
The logarithmic lower bound follows from
\[
\frac{d}{d\xi_{\mathrm{log}}}
\left\{\log(1+\xi_{\mathrm{log}})
-\frac{2\xi_{\mathrm{log}}}{\xi_{\mathrm{log}}+2}\right\}
=
\frac{\xi_{\mathrm{log}}^2}
{(1+\xi_{\mathrm{log}})(\xi_{\mathrm{log}}+2)^2}
\ge0,
\qquad \xi_{\mathrm{log}}\in[0,\infty),
\]
since the difference is zero at \(\xi_{\mathrm{log}}=0\). Applying it at \(\mathfrak x_{\mathrm{up}}/\lambda_\circ\) gives
\[
\vartheta_{\mathrm{up}}
\ge\frac{2\mathfrak x_{\mathrm{up}}}{2\lambda_\circ+\mathfrak x_{\mathrm{up}}},
\qquad
(\lambda_\circ+\mathfrak x_{\mathrm{up}})\vartheta_{\mathrm{up}}
-\mathfrak x_{\mathrm{up}}
\ge
\frac{\mathfrak x_{\mathrm{up}}^2}
{2\lambda_\circ+\mathfrak x_{\mathrm{up}}}
\ge z_{\mathrm{tail}}.
\]
The last inequality follows by expanding the threshold:
\[
\mathfrak x_{\mathrm{up}}^2
-z_{\mathrm{tail}}(2\lambda_\circ+\mathfrak x_{\mathrm{up}})
=
3z_{\mathrm{tail}}\sqrt{2\lambda_\circ z_{\mathrm{tail}}}
+2z_{\mathrm{tail}}^2\ge0.
\]
Thus the upper-tail probability is at most \(e^{-z_{\mathrm{tail}}}\). Taking \(z_{\mathrm{tail}}=40L\), and applying the lower-tail exponential Markov inequality with
\(e^{-t}-1+t\le t^2/2\), \(t\ge0\), gives
\[
\Pr\!\left\{
\mathsf N_\circ-\lambda_\circ>
\sqrt{80\lambda_\circ L}+80L
\right\}\le e^{-40L},
\qquad
\Pr\!\left\{
\lambda_\circ-\mathsf N_\circ>
\sqrt{80\lambda_\circ L}
\right\}\le e^{-40L}.
\]
The event
\[
\left\{
|\mathsf N_\circ-\lambda_\circ|
>\frac{1024}{4}\left(\sqrt{\mathsf N_\circ L}+L\right)
\right\}
\]
is contained in the union of these two tail events: on their complement, substitute
\(\lambda_\circ\le\mathsf N_\circ+|\mathsf N_\circ-\lambda_\circ|\) and square the nonnegative terms. Since \(H_0\ge1024\), failure of a pilot coordinate implies the corresponding event above. A union bound over the four coordinates proves the probability estimate.

For the moments, square-root subadditivity and
\(\sqrt{|\mathsf N_\circ-\lambda_\circ|L}
\le\{|\mathsf N_\circ-\lambda_\circ|+L\}/2\) imply
\[
\mathfrak e_\circ
\le1536\left(\mathfrak s_\circ+
|\mathsf N_\circ-\lambda_\circ|\right).
\]
Also \(\mathfrak s_\circ\ge1\),
\(\lambda_\circ\le\mathfrak s_\circ^2\), and
\(\lambda_\circ\vartheta_\circ^2\le1/4\). Using
\(|e^t-1-t|\le t^2\) for \(|t|\le1\) in the generating function therefore gives
\[
\mathbb E\exp\left(\frac{\mathfrak e_\circ}{4096\mathfrak s_\circ}\right)
\le
e^{1/2}\left(
\mathbb E e^{\vartheta_\circ(\mathsf N_\circ-\lambda_\circ)}
+
\mathbb E e^{-\vartheta_\circ(\mathsf N_\circ-\lambda_\circ)}
\right)
\le2e^{3/4}\le8.
\]
Thus, for \(j=2,4\), the inequality \(v^j\le j!e^v\), \(v\ge0\), yields
\[
\mathbb E\mathfrak e_\circ^j
\le8j!\,4096^j\mathfrak s_\circ^j
\le10^{4j+4}\mathfrak s_\circ^j.
\]
Substituting \(\lambda_\circ=mq_{ay,x}\) gives the displayed fourth-moment bound, because
\[
e_{ay,x}^{\mathrm p}
\le\frac{H_0}{1024m}\mathfrak e_\circ.
\]
We next localize the scalar moments before summing over failing coordinates. For the scalar Poisson count above, define
\[
\Delta_\circ=|\mathsf N_\circ-\lambda_\circ|,
\qquad
\mathcal B_\circ=
\left\{
\Delta_\circ>256\left(\sqrt{\mathsf N_\circ L}+L\right)
\right\}.
\]
Applying the exponential Markov calculations at level \(80L\) gives
\[
\Pr\!\left\{
\mathsf N_\circ-\lambda_\circ>
\sqrt{160\lambda_\circ L}+160L
\right\}\le e^{-80L},
\qquad
\Pr\!\left\{
\lambda_\circ-\mathsf N_\circ>
\sqrt{160\lambda_\circ L}
\right\}\le e^{-80L}.
\]
To compare \(\mathcal B_\circ\) with these events, first suppose \(\mathsf N_\circ\ge\lambda_\circ\). On \(\mathcal B_\circ\),
\[
\sqrt{160\lambda_\circ L}+160L
\le16\sqrt{\mathsf N_\circ L}+160L
\le256\left(\sqrt{\mathsf N_\circ L}+L\right)
<\Delta_\circ.
\]
If \(\mathsf N_\circ<\lambda_\circ\), then \(\lambda_\circ=\mathsf N_\circ+\Delta_\circ\). The defining inequality of \(\mathcal B_\circ\), and monotonicity of the quadratic above that threshold, give
\[
\begin{aligned}
\Delta_\circ^2-160L\Delta_\circ-160\mathsf N_\circ L
>{}&65376\mathsf N_\circ L
+90112L\sqrt{\mathsf N_\circ L}+24576L^2>0.
\end{aligned}
\]
Thus \(\Delta_\circ>\sqrt{160\lambda_\circ L}\). The two cases prove
\[
\Pr(\mathcal B_\circ)\le2e^{-80L}.
\]
The exponential score bound above also yields, for every integer \(1\le j_{\mathrm{mom}}\le8\),
\[
\mathbb E\mathfrak e_\circ^{j_{\mathrm{mom}}}
\le8j_{\mathrm{mom}}!\,4096^{j_{\mathrm{mom}}}
\mathfrak s_\circ^{j_{\mathrm{mom}}}
\le10^{4j_{\mathrm{mom}}+4}\mathfrak s_\circ^{j_{\mathrm{mom}}}.
\]
For \(j_{\mathrm{mom}}\in\{1,2,4\}\), define the positive balancing coefficient by
\[
\upsilon_{\circ,j_{\mathrm{mom}}}
=\frac{e^{-40L}}{\mathfrak s_\circ^{j_{\mathrm{mom}}}}.
\]
The nonnegative square
\((\upsilon_{\circ,j_{\mathrm{mom}}}\mathfrak e_\circ^{j_{\mathrm{mom}}}-1)^2\), on \(\mathcal B_\circ\), proves the pointwise inequality
\[
\mathfrak e_\circ^{j_{\mathrm{mom}}}\mathbf1_{\mathcal B_\circ}
\le\frac12\left(
\upsilon_{\circ,j_{\mathrm{mom}}}\mathfrak e_\circ^{2j_{\mathrm{mom}}}
+\upsilon_{\circ,j_{\mathrm{mom}}}^{-1}\mathbf1_{\mathcal B_\circ}
\right).
\]
Outside \(\mathcal B_\circ\), its left side is zero and its right side is nonnegative. Integration therefore gives the localized scalar bound
\[
\begin{aligned}
\mathbb E\!\left[
\mathfrak e_\circ^{j_{\mathrm{mom}}}\mathbf1_{\mathcal B_\circ}
\right]
&\le\frac{10^{8j_{\mathrm{mom}}+4}+2}{2}
 e^{-40L}\mathfrak s_\circ^{j_{\mathrm{mom}}}\\
&\le10^{8j_{\mathrm{mom}}+8}
 e^{-40L}\mathfrak s_\circ^{j_{\mathrm{mom}}}.
\end{aligned}
\]
% lean: Causalean.Stat.Concentration.PoissonSelfNormalized.poisson_weighted_bad_moment

For each \((a,y)\in\{0,1\}^2\) and each cell \(x\), define the raw score, raw local scale, and raw failure event by
\[
\mathfrak e_{ay,x}
=|N'_{ay,x}-mq_{ay,x}|
+1024\left(\sqrt{N'_{ay,x}L}+L\right),
\qquad
\mathfrak s_{ay,x}=\sqrt{mq_{ay,x}L}+L,
\]
\[
\mathcal B_{ay,x}
=\left\{
|N'_{ay,x}-mq_{ay,x}|>
256\left(\sqrt{N'_{ay,x}L}+L\right)
\right\},
\qquad
\mathcal B_x^{\cup}=\bigcup_{a,y}\mathcal B_{ay,x}.
\]
The interval-failure criterion and \(H_0\ge1024\) imply
\[
\mathcal G_x^c\subseteq\mathcal B_x^{\cup},
\qquad
e_{ay,x}^{\mathrm p}
\le\frac{H_0}{1024m}\mathfrak e_{ay,x}.
\]
For \(j_{\mathrm{mom}}\in\{1,2,4\}\), consider the contribution of one coordinate's score on another coordinate's failure event. If \((a,y)=(a',y')\), the localized scalar bound gives
\[
\mathbb E_{\mathrm{aux}}\!\left[
\mathfrak e_{ay,x}^{j_{\mathrm{mom}}}\mathbf1_{\mathcal B_{ay,x}}
\right]
\le10^{8j_{\mathrm{mom}}+8}e^{-40L}
\mathfrak s_{ay,x}^{j_{\mathrm{mom}}}.
\]
If \((a,y)\ne(a',y')\), independence of the pilot coordinates factors the expectation:
\[
\begin{aligned}
\mathbb E_{\mathrm{aux}}\!\left[
\mathfrak e_{ay,x}^{j_{\mathrm{mom}}}\mathbf1_{\mathcal B_{a'y',x}}
\right]
&=\mathbb E_{\mathrm{aux}}\mathfrak e_{ay,x}^{j_{\mathrm{mom}}}
\Pr(\mathcal B_{a'y',x})\\
&\le2\cdot10^{4j_{\mathrm{mom}}+4}e^{-40L}
\mathfrak s_{ay,x}^{j_{\mathrm{mom}}},
\end{aligned}
\]
where the scalar probability bound was enlarged from \(2e^{-80L}\) to \(2e^{-40L}\). Hence both cases satisfy
\[
\mathbb E_{\mathrm{aux}}\!\left[
\mathfrak e_{ay,x}^{j_{\mathrm{mom}}}\mathbf1_{\mathcal B_{a'y',x}}
\right]
\le
\left(10^{8j_{\mathrm{mom}}+8}+2\cdot10^{4j_{\mathrm{mom}}+4}\right)
e^{-40L}\mathfrak s_{ay,x}^{j_{\mathrm{mom}}}.
\]
The union indicator and the finite-sum power inequality give the pointwise bound
\[
\left(\sum_{a,y}\mathfrak e_{ay,x}\right)^{j_{\mathrm{mom}}}
\mathbf1_{\mathcal B_x^{\cup}}
\le4^{j_{\mathrm{mom}}-1}
\sum_{a',y'}\sum_{a,y}
\mathfrak e_{ay,x}^{j_{\mathrm{mom}}}
\mathbf1_{\mathcal B_{a'y',x}}.
\]
To sum the local scales, Cauchy--Schwarz and \(\sum_{a,y}q_{ay,x}=p_x\) give
\[
\sum_{a,y}\mathfrak s_{ay,x}
\le2\sqrt{mp_xL}+4L
\le4\left(\sqrt{mp_xL}+L\right),
\qquad
\sum_{a,y}\mathfrak s_{ay,x}^{j_{\mathrm{mom}}}
\le4\left(\sum_{a,y}\mathfrak s_{ay,x}\right)^{j_{\mathrm{mom}}}.
\]
Integrating the pointwise bound and inserting the two coordinate cases therefore yields
\[
\begin{aligned}
\mathbb E_{\mathrm{aux}}\!\left[
\left(\sum_{a,y}\mathfrak e_{ay,x}\right)^{j_{\mathrm{mom}}}
\mathbf1_{\mathcal B_x^{\cup}}
\right]
\le{}&4^{2j_{\mathrm{mom}}+1}
\left(10^{8j_{\mathrm{mom}}+8}+2\cdot10^{4j_{\mathrm{mom}}+4}\right)
e^{-40L}\left(\sqrt{mp_xL}+L\right)^{j_{\mathrm{mom}}}\\
\le{}&10^{80(j_{\mathrm{mom}}+1)}e^{-20L}
\left(\sqrt{mp_xL}+L\right)^{j_{\mathrm{mom}}}.
\end{aligned}
\]
The last enlargement uses, for the three moment orders,
\[
4^{2j_{\mathrm{mom}}+1}
\left(10^{8j_{\mathrm{mom}}+8}+2\cdot10^{4j_{\mathrm{mom}}+4}\right)
\le10^{80(j_{\mathrm{mom}}+1)},
\qquad e^{-40L}\le e^{-20L}.
\]
Finally, the raw-score comparison and
\[
\frac{\sqrt{mp_xL}+L}{m}
=\sqrt{p_x\delta_{\mathrm p}}+\delta_{\mathrm p}
\]
prove the normalized aggregate bound
\[
\mathbb E_{\mathrm{aux}}\!\left[
(e_x^{\mathrm p})^{j_{\mathrm{mom}}}\mathbf1_{\mathcal G_x^c}
\right]
\le
\left(\frac{H_0}{1024}\right)^{j_{\mathrm{mom}}}
10^{80(j_{\mathrm{mom}}+1)}e^{-20L}
\left(\sqrt{p_x\delta_{\mathrm p}}+\delta_{\mathrm p}\right)^{j_{\mathrm{mom}}}.
\]
Taking \(j_{\mathrm{mom}}=1\) gives the stated aggregate first-moment estimate with numerical constant \(10^{160}\).
% lean: CausalSmith.Stat.OptvalueVanishingoverlapRate.pilotRectangle_bad_probability_le
% lean: CausalSmith.Stat.OptvalueVanishingoverlapRate.pilotRectangle_bad_aggregate_moment_le

The midpoint error satisfies
\[
|b_{ay,x}-q_{ay,x}|\le e_{ay,x}^{\mathrm p},
\qquad
s(b_x)\le p_x+e_x^{\mathrm p}.
\]
Using the clipping-scale square bound from the preceding step, we obtain
\[
\begin{aligned}
\mathbb E_{\mathrm{aux}}
[(S_x^{\mathrm{aw}})^2\mathbf1_{\mathcal G_x^c}]
\le512H_0^2\bigg\{
8e^{-40L}\mathfrak a_x
+\frac{\delta_{\mathrm p}}{\epsilon}
\frac{H_0}{1024}10^{160}e^{-20L}
\left(\sqrt{p_x\delta_{\mathrm p}}+\delta_{\mathrm p}\right)
\bigg\}.
\end{aligned}
\]
Since
\[
\frac{\delta_{\mathrm p}}{\epsilon}
\left(\sqrt{p_x\delta_{\mathrm p}}+\delta_{\mathrm p}\right)
\le2\mathfrak a_x,
\]
where \(\sqrt{p_x\delta_{\mathrm p}}\le p_x+\delta_{\mathrm p}\) and \(\epsilon\le1\), this becomes
\[
\mathbb E_{\mathrm{aux}}
[(S_x^{\mathrm{aw}})^2\mathbf1_{\mathcal G_x^c}]
\le
C_{\mathrm{fail,s}}e^{-20L}\mathfrak a_x,
\]
with
\[
C_{\mathrm{fail,s}}
=
512H_0^2\left(8+2\frac{H_0}{1024}10^{160}\right).
\]

For an occupied cell, \cref{obj:lem:armwise-modulus} gives
\[
|F_\epsilon(b_x)-F_\epsilon(q_x)|
\le
4\sum_{a,y}\frac{p_x}{s_a(q_x)}e_{ay,x}^{\mathrm p}.
\]
The fourth-moment and probability estimates imply
\[
\begin{aligned}
\mathbb E_{\mathrm{aux}}
[(e_{ay,x}^{\mathrm p})^2\mathbf1_{\mathcal G_x^c}]
&\le
\left\{\mathbb E_{\mathrm{aux}}(e_{ay,x}^{\mathrm p})^4\right\}^{1/2}
\Pr(\mathcal G_x^c)^{1/2}\\
&\le
\sqrt{8\cdot10^{20}}
\left(\frac{H_0}{1024}\right)^2
e^{-20L}
\left(\sqrt{q_{ay,x}\delta_{\mathrm p}}+\delta_{\mathrm p}\right)^2.
\end{aligned}
\]
For each arm, the overlap floor gives
\[
\left(\frac{p_x}{s_a(q_x)}\right)^2
\sum_y
\left(\sqrt{q_{ay,x}\delta_{\mathrm p}}+\delta_{\mathrm p}\right)^2
\le8\mathfrak a_x.
\]
Indeed, the left side is at most
\[
2\frac{p_x^2\delta_{\mathrm p}}{s_a(q_x)}
+
4\frac{p_x^2\delta_{\mathrm p}^2}{s_a(q_x)^2}
\le
2\frac{p_x\delta_{\mathrm p}}{\epsilon}
+
4\frac{\delta_{\mathrm p}^2}{\epsilon^2}.
\]
Squaring the four-coordinate sum consequently yields
\[
\mathbb E_{\mathrm{aux}}
[(F_\epsilon(b_x)-F_\epsilon(q_x))^2\mathbf1_{\mathcal G_x^c}]
\le C_{\mathrm{fail,f}}e^{-20L}\mathfrak a_x,
\]
where
\[
C_{\mathrm{fail,f}}
=
1024\sqrt{8\cdot10^{20}}
\left(\frac{H_0}{1024}\right)^2.
\]

Clipping supplies the pathwise bound
\[
|T_x-F_\epsilon(q_x)|
\le
d^{1/4}S_x^{\mathrm{aw}}
+
|F_\epsilon(b_x)-F_\epsilon(q_x)|.
\]
Hence, with
\[
C_{\mathrm d}=2(C_{\mathrm{fail,s}}+C_{\mathrm{fail,f}}),
\]
we have
\[
\mathbb E_{\mathrm{aux}}
[(T_x-F_\epsilon(q_x))^2\mathbf1_{\mathcal G_x^c}]
\le C_{\mathrm d}d^{-4}\mathfrak a_x.
\]
Here
\[
d^{1/2}e^{-20L}\le d^{-4},
\qquad e^{-20L}\le d^{-4}
\]
for every \(d\ge1\). For null cells, \(\mathcal G_x^c\) has probability zero, so the same bound follows directly.
% lean: CausalSmith.Stat.OptvalueVanishingoverlapRate.markedPoissonCellValue_bad_sq_le

\textit{9. Independent-cell risk assembly.}
Define, on the whole auxiliary sample space,
\[
E_x^{\mathrm{good}}
=
(T_x-F_\epsilon(q_x))\mathbf1_{\mathcal G_x},
\qquad
E_x^{\mathrm{bad}}
=
(T_x-F_\epsilon(q_x))\mathbf1_{\mathcal G_x^c}.
\]
Then
\[
T_x-F_\epsilon(q_x)=E_x^{\mathrm{good}}+E_x^{\mathrm{bad}},
\qquad
E_x^{\mathrm{good}}E_x^{\mathrm{bad}}=0.
\]
All these variables have finite second moments by the preceding estimates. Independence of the complete cell statistics gives
\[
\mathbb E_{\mathrm{aux}}
\left(\sum_xT_x-\Psi(\mathbb P)\right)^2
=
\sum_x\operatorname{Var}_{\mathrm{aux}}(T_x)
+
\left\{\sum_x
\mathbb E_{\mathrm{aux}}(T_x-F_\epsilon(q_x))\right\}^2.
\]
The disjoint error decomposition implies
\[
\operatorname{Var}_{\mathrm{aux}}(T_x)
\le
\mathbb E_{\mathrm{aux}}(E_x^{\mathrm{good}})^2
+
\mathbb E_{\mathrm{aux}}(E_x^{\mathrm{bad}})^2.
\]
Applying \((v+w)^2\le2v^2+2w^2\) to the total bias, and Cauchy--Schwarz to its failed-pilot part, therefore gives
\[
\begin{aligned}
\mathbb E_{\mathrm{aux}}
\left(\sum_xT_x-\Psi(\mathbb P)\right)^2
\le{}&
\sum_x\mathbb E_{\mathrm{aux}}(E_x^{\mathrm{good}})^2
+
2\left\{\sum_x\mathbb E_{\mathrm{aux}}E_x^{\mathrm{good}}\right\}^2\\
&+
(1+2d)\sum_x\mathbb E_{\mathrm{aux}}(E_x^{\mathrm{bad}})^2.
\end{aligned}
\]
% lean: CausalSmith.Stat.OptvalueVanishingoverlapRate.independentCell_good_bad_sqRisk_le

\textit{10. Absorption of the pilot scales.}
Assume the active-branch condition \(d/L\le n\epsilon\), and define
\[
\mathfrak r_{\mathrm{aux}}=\frac{d}{m\epsilon L},
\qquad
\mathfrak m_{\mathrm{sum}}
=\sum_x\mathfrak a_x
=\frac L{m\epsilon}+\frac{dL^2}{m^2\epsilon^2}.
\]
Then
\[
0<\mathfrak r_{\mathrm{aux}}\le8,
\qquad
\mathfrak m_{\mathrm{sum}}
\le\frac{L^2+8L^4}{m\epsilon L}.
\]
The second inequality follows by substituting
\(d\le8m\epsilon L\) into its second summand.

For a positive real \(\delta_{\mathrm{dec}}\), define the numerical envelope
\[
\mathcal M(\delta_{\mathrm{dec}})
=
\left(\frac{e^{\delta_{\mathrm{dec}}/2}}{\delta_{\mathrm{dec}}/2}\right)^2
+
8\left(\frac{e^{\delta_{\mathrm{dec}}/4}}{\delta_{\mathrm{dec}}/4}\right)^4.
\]
The elementary inequality \(\log z\le z^u/u\), \(z\ge1\), \(u>0\), applied with \(z=ed\), gives
\[
(L^2+8L^4)d^{-\delta_{\mathrm{dec}}}
\le\mathcal M(\delta_{\mathrm{dec}}).
\]
Set
\[
C_{\mathrm{var}}=\mathcal M(15/16),\qquad
C_{\mathrm{bias}}=\mathcal M(3/8),\qquad
C_{\mathrm{bad,sum}}=\mathcal M(1).
\]
Thus
\[
d^{1/16}\sum_x\mathbb E_{\mathrm{aux}}(S_x^{\mathrm{aw}})^2
\le C_{\mathrm s}C_{\mathrm{var}}\mathfrak r_{\mathrm{aux}}.
\]

The probability-mass identity and Cauchy--Schwarz give
\[
\sum_x\eta_x\le\sqrt{\mathfrak r_{\mathrm{aux}}}+\mathfrak r_{\mathrm{aux}},
\qquad
\sum_x\eta_x^2\le\left(\sum_x\eta_x\right)^2
\le18\mathfrak r_{\mathrm{aux}}.
\]
Also,
\[
\left(\sum_x\mathbb E_{\mathrm{aux}}S_x^{\mathrm{aw}}\right)^2
\le d\sum_x\mathbb E_{\mathrm{aux}}(S_x^{\mathrm{aw}})^2,
\]
and hence
\[
d^{-3/8}
\left(\sum_x\mathbb E_{\mathrm{aux}}S_x^{\mathrm{aw}}\right)^2
\le C_{\mathrm s}C_{\mathrm{bias}}\mathfrak r_{\mathrm{aux}}.
\]
The successful-pilot estimates now imply
\[
\sum_x\mathbb E_{\mathrm{aux}}(E_x^{\mathrm{good}})^2
\le
(C_{\mathrm v}C_{\mathrm s}C_{\mathrm{var}}+18C_{\mathrm e})
\mathfrak r_{\mathrm{aux}},
\]
\[
\left(\sum_x\mathbb E_{\mathrm{aux}}E_x^{\mathrm{good}}\right)^2
\le
2C_{\mathrm b}^2(18+C_{\mathrm s}C_{\mathrm{bias}})
\mathfrak r_{\mathrm{aux}}.
\]
For the failed-pilot term, \(d\,d^{-4}\le1\) gives
\[
d\sum_x\mathbb E_{\mathrm{aux}}(E_x^{\mathrm{bad}})^2
\le C_{\mathrm d}\mathfrak m_{\mathrm{sum}}
\le C_{\mathrm d}C_{\mathrm{bad,sum}}\mathfrak r_{\mathrm{aux}},
\]
so, since \(d\ge1\),
\[
(1+2d)\sum_x\mathbb E_{\mathrm{aux}}(E_x^{\mathrm{bad}})^2
\le3C_{\mathrm d}C_{\mathrm{bad,sum}}\mathfrak r_{\mathrm{aux}}.
\]

Define
\[
\begin{aligned}
C_{\mathrm P}={}&
C_{\mathrm v}C_{\mathrm s}C_{\mathrm{var}}
+4C_{\mathrm b}^2(18+C_{\mathrm s}C_{\mathrm{bias}})
+3C_{\mathrm d}C_{\mathrm{bad,sum}}
+18C_{\mathrm e}+1.
\end{aligned}
\]
This is a positive finite universal constant. The independent-cell assembly yields
\[
\mathbb E_{\mathrm{aux}}
\left(\sum_xT_x-\Psi(\mathbb P)\right)^2
\le C_{\mathrm P}\frac{d}{m\epsilon L}.
\]
Finally, define the projected auxiliary statistic by
\[
V_{\mathrm{aux}}=\Pi_{[0,1]}\left(\sum_xT_x\right).
\]
Since \(\Psi(\mathbb P)\in[0,1]\), projection contracts its squared error, giving
\[
\mathbb E_{\mathrm{aux}}
(V_{\mathrm{aux}}-\Psi(\mathbb P))^2
\le C_{\mathrm P}\frac{d}{m\epsilon L}.
\]
% lean: CausalSmith.Stat.OptvalueVanishingoverlapRate.activeBranch_good_bad_approximation_sqRisk_rate_le
% lean: CausalSmith.Stat.OptvalueVanishingoverlapRate.markedPoissonStatistic_goodPilotEstimates_rate_le
% lean: CausalSmith.Stat.OptvalueVanishingoverlapRate.markedPoissonStatistic_active_rate_tuning_exists

\textit{11. Fixed-sample transfer.}
Set
\[
D_0=2,\qquad
C_\star=3+8\log(2e)+8C_{\mathrm P}.
\]
Then \(D_0\ge2\) and \(0<C_\star<\infty\). For the admissible theorem parameters \(n\ge1,d\ge2\), the finite sample alphabet makes the estimator in \cref{obj:def:armwise-estimator} measurable. Its auxiliary variables are integrated by the specified finite sum, so it is a deterministic function of the observations.

Fix an admissible observed law and its sampling law from \cref{obj:ass:iid-sampling}. First consider \(d<D_0\). With the chosen cutoff this branch has an empty parameter range, because \(d\ge2=D_0\).

Next suppose \(d\ge D_0\) and \(n\epsilon<d/L\). The estimator equals \(1/2\), and
\[
\mathbb E_{\mathbb P}
\left(\widehat V_{n,d,\epsilon}^{\mathrm{AF}}-\Psi(\mathbb P)\right)^2
\le\frac14,
\qquad
\min\left\{1,\frac{d}{n\epsilon L}\right\}=1.
\]
This proves the desired bound on that branch.

On the active branch, put the auxiliary count and marks alongside the fixed observed sample. Conditional averaging of squared error gives
\[
\begin{aligned}
\mathbb E_{\mathbb P}
\left(\widehat V_{n,d,\epsilon}^{\mathrm{AF}}-\Psi(\mathbb P)\right)^2
\le{}&
\mathbb E_{\mathrm{aux}}
\left[
(V_{\mathrm{aux}}-\Psi(\mathbb P))^2
\mathbf1_{\{\mathsf M\le n\}}
\right]\\
&+\Pr(\mathsf M>n)\Psi(\mathbb P)^2.
\end{aligned}
\]
Indeed, the auxiliary mixture has value zero on overflow. For each \(k\le n\), the first \(k\) observations have law \(\mathbb P^{\otimes k}\); consequently the integrated nonoverflow loss is exactly the corresponding part of the unrestricted finite-Poisson experiment.

The quarter-mean Poisson cap satisfies
\[
\Pr(\mathsf M>n)\le e^{-n/2}.
\]
For example, exponential Markov inequality at \(t=\log4\) gives
\[
\Pr(\mathsf M>n)
\le e^{-n(\log4-3/4)}
\le e^{-n/2}.
\]
Also, \(e^{n/2}\ge n/2\), \(L\le d\), and \(\epsilon\le1/2\) imply
\[
e^{-n/2}\le\frac2n
\le2\frac{d}{n\epsilon L}.
\]
Since \(\Psi(\mathbb P)^2\le1\), the auxiliary risk bound and \(m=n/8\) therefore give
\[
\mathbb E_{\mathbb P}
\left(\widehat V_{n,d,\epsilon}^{\mathrm{AF}}-\Psi(\mathbb P)\right)^2
\le
(8C_{\mathrm P}+2)\frac{d}{n\epsilon L}.
\]
The estimator and target both lie in \([0,1]\), so the risk is also at most one. Splitting according as \(d/(n\epsilon L)\ge1\) or \(d/(n\epsilon L)\le1\), we obtain
\[
\mathbb E_{\mathbb P}
\left(\widehat V_{n,d,\epsilon}^{\mathrm{AF}}-\Psi(\mathbb P)\right)^2
\le
(8C_{\mathrm P}+2)
\min\left\{1,\frac{d}{n\epsilon L}\right\}
\le
C_\star\min\left\{1,\frac{d}{n\epsilon L}\right\}.
\]
% lean: CausalSmith.Stat.OptvalueVanishingoverlapRate.armwiseEstimator_active_le_markedPoisson_sqRisk
% lean: CausalSmith.Stat.OptvalueVanishingoverlapRate.armwiseEstimator_allBranches_rate_le

\textit{12. Uniformity over sampling laws.}
The preceding estimates hold for every admissible observed law and every family of sampling laws satisfying \cref{obj:ass:iid-sampling}. If the admissible class is empty, its worst-risk value, with the empty-class convention zero, satisfies the bound because the right side is nonnegative. Otherwise, taking the supremum of the pointwise bound gives
\[
\sup_{\mathbb P\in\mathcal D_{d,\epsilon}^{\mathrm{obs}}}
\mathbb E_{\mathbb P}
\left[
\left\{
\widehat V_{n,d,\epsilon}^{\mathrm{AF}}-\Psi(\mathbb P)
\right\}^{2}
\right]
\le
C_\star\min\left\{1,\frac{d}{n\epsilon L_d}\right\}.
\]
Together with measurability and determinism, this proves the claim.
% lean: CausalSmith.Stat.OptvalueVanishingoverlapRate.armwiseEstimator_allBranches_worstRisk_le
% lean: CausalSmith.Stat.OptvalueVanishingoverlapRate.observable_upper
\end{proof}

\begin{proof}[Proof of \cref{obj:thm:weighted-separation-and-frontier}]
\begin{enumerate}
\item
\textbf{Approximation and constrained priors.}\quad
We first establish the packet lower bound. Throughout the approximation argument, assume
\[
K\ge2,\qquad 2\le M\le K^2,\qquad 0\le\epsilon\le\tfrac12,
\]
and write
\[
\mathcal E=\mathcal E_{K,M,\epsilon}^{w},
\qquad
\Delta=\Delta_{K,M,\epsilon},
\qquad
z_M(\omega)=\frac M2(1+\cos\omega),
\qquad
\vartheta_M=\arccos(2/M-1).
\]
Thus \(z_M(\vartheta_M)=1\). For integers \(N\ge2\), use
\(\mathcal J_N,k_N,G_N\) from
\cref{obj:lem:jackson-packet-localization}, and define the Lebesgue-normalized kernel
\[
J_N^{\mathrm J}(\omega)=\frac{\mathcal J_N(\omega)}{2\pi},
\qquad \omega\in\mathbb R.
\]
We will use the following two kernel moments:
\[
\int_{-\pi}^{\pi}\omega^2J_N^{\mathrm J}(\omega)\,d\omega
\le\frac{64}{N^2},
\qquad
\int_{-\pi}^{\pi}|\omega|J_N^{\mathrm J}(\omega)\,d\omega
\le\frac{32}{N}.
\]
Here is a direct derivation. Set, with removable values filled continuously,
\[
D_N(\omega)=\frac{\sin(N\omega/2)}{\sin(\omega/2)},
\qquad
\mathcal Z_N^{\mathrm J}
=\int_{-\pi}^{\pi}D_N(\omega)^4\,d\omega.
\]
The finite geometric-sum identity and orthogonality give
\[
D_N(\omega)^2
=N+2\sum_{\ell=1}^{N-1}(N-\ell)\cos(\ell\omega),
\]
\[
\int_{-\pi}^{\pi}D_N(\omega)^2\,d\omega=2\pi N,
\qquad
\mathcal Z_N^{\mathrm J}
=\frac{2\pi}{3}N(2N^2+1)
\ge\frac{4\pi}{3}N^3,
\qquad
J_N^{\mathrm J}=\frac{D_N^4}{\mathcal Z_N^{\mathrm J}}.
\]
For \(|\omega|\le\pi\), the inequality
\(|\sin(\omega/2)|\ge|\omega|/\pi\) implies
\[
\omega^2D_N(\omega)^4\le\pi^2D_N(\omega)^2.
\]
Consequently,
\[
\int_{-\pi}^{\pi}\omega^2J_N^{\mathrm J}(\omega)\,d\omega
\le
\frac{2\pi^3N}{(4\pi/3)N^3}
\le\frac{24}{N^2}
\le\frac{64}{N^2}.
\]
Finally, integrating
\[
|\omega|\le\frac N4\omega^2+\frac1N
\]
against \(J_N^{\mathrm J}\), whose integral is one, proves
\[
\int_{-\pi}^{\pi}|\omega|J_N^{\mathrm J}(\omega)\,d\omega
\le\frac N4\frac{24}{N^2}+\frac1N
\le\frac{32}{N}.
\]
% lean: Causalean.Mathlib.Analysis.JacksonApproximation.jackson_first_moment

\item
Define the nonlinear remainder and its angular version by
\[
g_\epsilon(z)
=\epsilon(1-z)_+
 +(1-2\epsilon)\frac{(z-1)_+}{1+z},
\qquad z\ge0,
\qquad
f_{M,\epsilon}(\omega)=g_\epsilon\{z_M(\omega)\},
\qquad \omega\in\mathbb R.
\]
Splitting at \(z=1\) gives
\[
\phi_\epsilon(z)=\epsilon(z-1)+g_\epsilon(z).
\]
For later integration, define the smooth curvature on the entire circle, including the kink locations, by
\[
a_{M,\epsilon}(\omega)=
\begin{cases}
\displaystyle \epsilon\frac M2\cos\omega,
&z_M(\omega)<1,\\[4pt]
\displaystyle
-\frac{4(1-2\epsilon)\{z_M'(\omega)\}^2}
       {\{1+z_M(\omega)\}^3}
+\frac{2(1-2\epsilon)z_M''(\omega)}
       {\{1+z_M(\omega)\}^2},
&z_M(\omega)>1,\\[6pt]
0,&z_M(\omega)=1.
\end{cases}
\]
On each open smooth arc this equals \(f_{M,\epsilon}''\). Since
\[
z_M'(\omega)=-\frac M2\sin\omega,
\qquad
z_M''(\omega)=-\frac M2\cos\omega,
\qquad
\{z_M'(\omega)\}^2=z_M(\omega)\{M-z_M(\omega)\},
\]
we have
\[
|a_{M,\epsilon}(\omega)|\le5M.
\]
Indeed, the lower branch is bounded by \(M/4\), while on the upper branch
\[
|a_{M,\epsilon}(\omega)|
\le
\frac{4Mz_M(\omega)}{\{1+z_M(\omega)\}^3}
+\frac{M}{\{1+z_M(\omega)\}^2}
\le5M.
\]
The scalar derivative jump is
\[
g_\epsilon'(1+)-g_\epsilon'(1-)
=\frac{1-2\epsilon}{2}+\epsilon=\frac12.
\]
Thus the angular derivative has a positive jump of size
\[
J_M^{\mathrm{jump}}=\frac{\sqrt{M-1}}2
\]
at each of \(-\vartheta_M\) and \(\vartheta_M\).

Twice integrating by parts on the three arcs
\([-\pi,-\vartheta_M]\), \([-\vartheta_M,\vartheta_M]\), and
\([\vartheta_M,\pi]\) therefore yields
\[
\begin{aligned}
\int_{-\pi}^{\pi}
f_{M,\epsilon}(\omega)k_N(\omega-\vartheta_M)\,\frac{d\omega}{2\pi}
={}&
\int_{-\pi}^{\pi}
a_{M,\epsilon}(\omega)G_N(\omega-\vartheta_M)\,
\frac{d\omega}{2\pi}\\
&+\frac{\sqrt{M-1}}{4\pi}
 \{G_N(0)+G_N(2\vartheta_M)\}.
\end{aligned}
\]
The endpoint terms at \(-\pi\) and \(\pi\) cancel by periodicity; the two remaining boundary terms are the displayed derivative jumps. The bounds in
\cref{obj:lem:jackson-packet-localization} supply universal constants
\(c_{\mathrm{loc}},C_{\mathrm{loc}}>0\) such that
\[
\frac{c_{\mathrm{loc}}}{N}\le|G_N(0)|,
\qquad
\int_{-\pi}^{\pi}|G_N(\omega)|\,\frac{d\omega}{2\pi}
\le\frac{C_{\mathrm{loc}}}{N^2}.
\]
In particular, with
\[
C_{\mathrm{sm}}=5C_{\mathrm{loc}},
\]
the smooth contribution satisfies
\[
\left|
\int_{-\pi}^{\pi}
a_{M,\epsilon}(\omega)G_N(\omega-\vartheta_M)\,
\frac{d\omega}{2\pi}
\right|
\le\frac{C_{\mathrm{sm}}M}{N^2}.
\]

To control the other kink, put
\[
\beta_M=\pi-\vartheta_M.
\]
Since \(M\ge2\), one has \(\vartheta_M\in[\pi/2,\pi)\), and hence
\[
\|2\vartheta_M\|_{\mathbb T}=2\beta_M,
\qquad
\cos\beta_M=1-\frac2M.
\]
The inequality \(1-\cos\beta_M\le\beta_M^2/2\) gives
\[
\beta_M\ge\frac2{\sqrt M},
\qquad
\|2\vartheta_M\|_{\mathbb T}\ge\frac4{\sqrt M}.
\]
Applying the localized bound for \(G_N\) gives, with
\[
C_{\mathrm{opp}}=\frac{C_{\mathrm{loc}}}{16},
\]
\[
|G_N(2\vartheta_M)|
\le
\frac{C_{\mathrm{loc}}}
     {N\{1+N\|2\vartheta_M\|_{\mathbb T}\}^2}
\le\frac{C_{\mathrm{opp}}M}{N^3}.
\]

Choose a universal integer \(A_0\) satisfying
\[
A_0>
\max\left\{
1,\frac{2C_{\mathrm{opp}}}{c_{\mathrm{loc}}},
\frac{32\pi C_{\mathrm{sm}}}{c_{\mathrm{loc}}}
\right\},
\qquad
N=A_0K.
\]
Then
\[
2C_{\mathrm{opp}}\le c_{\mathrm{loc}}A_0^2,
\qquad
32\pi C_{\mathrm{sm}}\le c_{\mathrm{loc}}A_0,
\]
so \(M\le K^2\) implies
\[
|G_N(2\vartheta_M)|\le\frac{c_{\mathrm{loc}}}{2N},
\qquad
|G_N(0)+G_N(2\vartheta_M)|
\ge\frac{c_{\mathrm{loc}}}{2N}.
\]
Also,
\[
\sqrt{M-1}\ge\frac{\sqrt M}{2},
\qquad
\frac{C_{\mathrm{sm}}M}{N^2}
\le\frac{c_{\mathrm{loc}}\sqrt M}{32\pi N}.
\]
The integration-by-parts identity consequently proves
\[
\left|
\int_{-\pi}^{\pi}
f_{M,\epsilon}(\omega)k_N(\omega-\vartheta_M)\,
\frac{d\omega}{2\pi}
\right|
\ge
\frac{c_{\mathrm{loc}}\sqrt M}{32\pi N}.
\]
% lean: CausalSmith.Stat.OptvalueVanishingoverlapRate.packetRemainder_response_lower

\item
Define the two-center packet and its normalization by
\[
H_{K,M}(\omega)
=k_N(\omega-\vartheta_M)+k_N(\omega+\vartheta_M),
\qquad
W_{K,M}
=\int_{-\pi}^{\pi}
\{1+z_M(\omega)\}|H_{K,M}(\omega)|\,\frac{d\omega}{2\pi}.
\]
The frequency exclusions in
\cref{obj:lem:jackson-packet-localization} imply
\[
\int_{-\pi}^{\pi}
P_{\mathrm{test}}\{z_M(\omega)\}H_{K,M}(\omega)\,
\frac{d\omega}{2\pi}=0,
\qquad
\deg P_{\mathrm{test}}\le K.
\]
Indeed, composing such a polynomial with \(z_M\) produces a trigonometric polynomial with frequencies of absolute value at most \(K\), whereas the translated packets have frequencies between \(2N+2\) and \(6N-2\). Translation preserves these exclusions.

Both \(f_{M,\epsilon}\) and \(k_N\) are even, so their pairings with the two translates agree. The affine term in
\(\phi_\epsilon=\epsilon(z-1)+g_\epsilon\) is annihilated. Thus
\[
\begin{aligned}
\left|
\int_{-\pi}^{\pi}
\phi_\epsilon\{z_M(\omega)\}H_{K,M}(\omega)\,
\frac{d\omega}{2\pi}
\right|
&=
2\left|
\int_{-\pi}^{\pi}
f_{M,\epsilon}(\omega)k_N(\omega-\vartheta_M)\,
\frac{d\omega}{2\pi}
\right|\\
&\ge
\frac{c_{\mathrm{loc}}}{16\pi A_0}\frac{\sqrt M}{K}.
\end{aligned}
\]

We next verify the normalization over the whole parameter range. The nonnegative Jackson kernel has a strict maximum at the periodic origin: the geometric-sum representation of \(D_N\) gives
\[
|D_N(\omega)|<N
\quad\text{when }\cos\omega<1,
\qquad
|D_N(0)|=N.
\]
Since \(\cos(2\vartheta_M)<1\),
\[
H_{K,M}(\vartheta_M)
=\mathcal J_N(0)
 +\mathcal J_N(2\vartheta_M)\cos(8N\vartheta_M)
\ge\mathcal J_N(0)-\mathcal J_N(2\vartheta_M)>0.
\]
Continuity and nonnegativity of the integrand defining \(W_{K,M}\) therefore give \(W_{K,M}>0\).

For the upper bound, use
\[
|H_{K,M}(\omega)|
\le
\mathcal J_N(\omega-\vartheta_M)
+\mathcal J_N(\omega+\vartheta_M)
\]
and the identity
\[
z_M(\omega+\vartheta_M)+z_M(\omega-\vartheta_M)
=2+(M-2)(1-\cos\omega).
\]
After translating the two periodic integrals,
\[
\begin{aligned}
W_{K,M}
&\le
\int_{-\pi}^{\pi}
\{4+(M-2)(1-\cos\omega)\}J_N^{\mathrm J}(\omega)\,d\omega\\
&\le
4+\frac{M-2}{2}
 \int_{-\pi}^{\pi}\omega^2J_N^{\mathrm J}(\omega)\,d\omega\\
&\le4+\frac{32(M-2)}{N^2}
\le36.
\end{aligned}
\]
Define
\[
c_{\mathrm p}=\frac{c_{\mathrm{loc}}}{576\pi A_0}>0.
\]
Dividing the packet response by \(W_{K,M}\le36\) proves
\[
\left|
\int_{-\pi}^{\pi}
\phi_\epsilon\{z_M(\omega)\}
\frac{H_{K,M}(\omega)}{W_{K,M}}\,
\frac{d\omega}{2\pi}
\right|
\ge c_{\mathrm p}\frac{\sqrt M}{K}.
\]
% lean: CausalSmith.Stat.OptvalueVanishingoverlapRate.packetNormalized_response_lower

\item
The packet parts in the statement are explicitly
\[
\sigma_+
=(z_M)_\#
\left[
\left(\frac{H_{K,M}}{W_{K,M}}\right)_+
\frac{d\omega}{2\pi}
\right],
\qquad
\sigma_-
=(z_M)_\#
\left[
\left(-\frac{H_{K,M}}{W_{K,M}}\right)_+
\frac{d\omega}{2\pi}
\right],
\]
where both measures inside the brackets are restricted to
\([-\pi,\pi]\). They are finite and supported on \([0,M]\), and therefore have integrable first moments. The packet is even, and the fibers of \(z_M\) consist of points with equal packet sign; these are the positive and negative Jordan parts of the pushed signed packet.

The annihilation and normalization established above give
\[
\int z^\ell\,d\sigma_+
=\int z^\ell\,d\sigma_-,
\qquad \ell=0,\ldots,K,
\]
\[
\int(1+z)\,d(\sigma_++\sigma_-)=1.
\]
Writing, as in \cref{obj:def:weighted-handle},
\[
t=\sigma_+(\mathbb R),
\qquad
u=\int z\,d\sigma_+,
\qquad
z_0=\frac{1-u}{1-t},
\]
the degree-zero and degree-one identities imply
\[
\sigma_-(\mathbb R)=t,
\qquad
\int z\,d\sigma_-=u,
\qquad
t+u=\frac12.
\]
Since \(t,u\ge0\),
\[
0\le t\le\frac12,
\qquad
z_0=\frac{1/2+t}{1-t},
\qquad
\frac12\le z_0\le2\le M.
\]
In particular \(t<1\), so the completion map is defined, and its positive-part coefficient equals \(1-t\). Set
\[
\nu_0=\sigma_-+(1-t)\delta_{z_0},
\qquad
\nu_1=\sigma_++(1-t)\delta_{z_0}.
\]
Their masses and first moments satisfy
\[
\nu_0(\mathbb R)=\nu_1(\mathbb R)=t+(1-t)=1,
\qquad
\int z\,d\nu_0=\int z\,d\nu_1=u+(1-t)z_0=1.
\]
The common atom cancels in every moment difference, so
\[
(\nu_0,\nu_1)\in\mathfrak N_{K,M},
\qquad
\mathscr H^w(\sigma_+,\sigma_-)=(\nu_0,\nu_1).
\]
It also cancels in the target difference:
\[
\begin{aligned}
\int\phi_\epsilon\,d\nu_1-\int\phi_\epsilon\,d\nu_0
&=
\int_{-\pi}^{\pi}
\phi_\epsilon\{z_M(\omega)\}
\frac{H_{K,M}(\omega)}{W_{K,M}}\,
\frac{d\omega}{2\pi}.
\end{aligned}
\]
Hence the explicit completion has gap at least
\(c_{\mathrm p}\sqrt M/K\).
% lean: CausalSmith.Stat.OptvalueVanishingoverlapRate.packetCompletion_separation_lower

\item
For any real polynomial \(P_{\mathrm{test}}\) of degree at most \(K\), define
\[
e(P_{\mathrm{test}})
=\sup_{z\in[0,M]}
\frac{|\phi_\epsilon(z)-P_{\mathrm{test}}(z)|}{1+z}.
\]
The supremum is finite by continuity. For any feasible prior pair, polynomial moment matching gives
\[
\int P_{\mathrm{test}}\,d\nu_1
=\int P_{\mathrm{test}}\,d\nu_0,
\]
and probability normalization with mean one gives
\[
\int(1+z)\,d\nu_0=\int(1+z)\,d\nu_1=2.
\]
Therefore
\[
\begin{aligned}
\left|\int\phi_\epsilon\,d\nu_1-\int\phi_\epsilon\,d\nu_0\right|
&=
\left|
\int(\phi_\epsilon-P_{\mathrm{test}})\,d\nu_1
-\int(\phi_\epsilon-P_{\mathrm{test}})\,d\nu_0
\right|\\
&\le
e(P_{\mathrm{test}})
\sum_{\iota=0}^1\int(1+z)\,d\nu_\iota
=4e(P_{\mathrm{test}}).
\end{aligned}
\]
Applying this bound to the packet completion and then taking the polynomial infimum yields, with
\[
c_{\mathrm E}=\frac{c_{\mathrm p}}4,
\]
\[
c_{\mathrm E}\frac{\sqrt M}{K}\le\mathcal E.
\]
In particular \(\mathcal E>0\).
% lean: CausalSmith.Stat.OptvalueVanishingoverlapRate.packetWeightedApproxError_lower

\item
We now prove the reverse comparison \(\mathcal E\le\Delta\) by a finite weighted dual construction. Define
\[
F_{\mathrm w}(z)=\frac{\phi_\epsilon(z)}{1+z},
\qquad
\mathcal V_K^{\mathrm w}
=\left\{
z\mapsto\frac{P_{\mathrm{test}}(z)}{1+z}:
\deg P_{\mathrm{test}}\le K
\right\}
\subset C([0,M],\mathbb R).
\]
This is a finite-dimensional subspace. Its closed ball of radius
\[
R_{\mathrm{best}}=2\|F_{\mathrm w}\|_\infty+1
\]
is compact, so the distance from \(F_{\mathrm w}\) attains a minimum there. The zero function gives distance at most \(\|F_{\mathrm w}\|_\infty\), while a function outside this ball has distance at least
\[
\|v_{\mathrm w}\|_\infty-\|F_{\mathrm w}\|_\infty
>\|F_{\mathrm w}\|_\infty+1.
\]
Thus a polynomial \(P_{\mathrm{best}}\) attains the global infimum:
\[
\deg P_{\mathrm{best}}\le K,
\qquad
e(P_{\mathrm{best}})=\mathcal E.
\]
Define its weighted residual and extremal set by
\[
r_{\mathrm{best}}(z)
=\frac{\phi_\epsilon(z)-P_{\mathrm{best}}(z)}{1+z},
\qquad
\mathcal A_{\mathrm{ext}}
=\{z\in[0,M]:|r_{\mathrm{best}}(z)|=\mathcal E\}.
\]

There exist ordered nodes and an orientation satisfying
\[
0\le z_0^{\mathrm{dual}}<\cdots<z_{K+1}^{\mathrm{dual}}\le M,
\qquad
o_{\mathrm{dual}}\in\{-1,1\},
\]
\[
r_{\mathrm{best}}(z_\ell^{\mathrm{dual}})
=o_{\mathrm{dual}}(-1)^\ell\mathcal E,
\qquad \ell=0,\ldots,K+1.
\]
To see this, suppose such nodes are absent. The positive and negative extremal sets are disjoint compact sets. If one is empty, the constant polynomial with the other sign agrees with the residual sign on all extrema. If both are nonempty, their positive separation permits a finite ordered cover by sufficiently small neighborhoods. Consecutive sign blocks in that cover have at most \(K\) transitions: \(K+1\) transitions would furnish the forbidden \(K+2\) alternating nodes. Placing one root in each transition gap and choosing the overall sign produces a polynomial \(Q_{\mathrm{sign}}\) such that
\[
\deg Q_{\mathrm{sign}}\le K,
\qquad
r_{\mathrm{best}}(z)Q_{\mathrm{sign}}(z)>0
\quad(z\in\mathcal A_{\mathrm{ext}}).
\]
Put
\[
v_{\mathrm{sign}}(z)=\frac{Q_{\mathrm{sign}}(z)}{1+z}.
\]
Define the set where the perturbation has the wrong sign or vanishes by
\[
\mathcal B_{\mathrm{bad}}
=\{z\in[0,M]:r_{\mathrm{best}}(z)v_{\mathrm{sign}}(z)\le0\}.
\]
This set is compact. If it is nonempty, the continuous absolute residual attains its maximum there, and
\[
\max_{z\in\mathcal B_{\mathrm{bad}}}|r_{\mathrm{best}}(z)|<\mathcal E.
\]
Indeed, equality at a maximizing point would make that point an extremum, where the residual and perturbation have strictly positive product. Define the threshold in both cases by
\[
b_{\mathrm{ext}}=
\begin{cases}
\displaystyle\frac{\max_{z\in\mathcal B_{\mathrm{bad}}}|r_{\mathrm{best}}(z)|+\mathcal E}{2},
&\mathcal B_{\mathrm{bad}}\ne\varnothing,\\[6pt]
\mathcal E/2,&\mathcal B_{\mathrm{bad}}=\varnothing.
\end{cases}
\]
Thus \(b_{\mathrm{ext}}<\mathcal E\), and in either case
\[
|r_{\mathrm{best}}(z)|\ge b_{\mathrm{ext}}
\quad\Longrightarrow\quad
r_{\mathrm{best}}(z)v_{\mathrm{sign}}(z)>0.
\]
The high-residual set
\[
\mathcal A_{\mathrm{high}}
=\{z\in[0,M]:|r_{\mathrm{best}}(z)|\ge b_{\mathrm{ext}}\}
\]
is compact and nonempty, since the residual attains its norm \(\mathcal E\). Consequently the following extrema exist:
\[
m_{\mathrm{ext}}
=\min_{z\in\mathcal A_{\mathrm{high}}}
 r_{\mathrm{best}}(z)v_{\mathrm{sign}}(z)>0,
\qquad
B_{\mathrm{sign}}
=\max_{z\in[0,M]}|v_{\mathrm{sign}}(z)|>0.
\]
The second constant is positive because the perturbation is nonzero at every residual extremum. These definitions give
\[
r_{\mathrm{best}}(z)v_{\mathrm{sign}}(z)\ge m_{\mathrm{ext}}
\quad\text{if }|r_{\mathrm{best}}(z)|\ge b_{\mathrm{ext}},
\qquad
|v_{\mathrm{sign}}(z)|\le B_{\mathrm{sign}}
\quad(z\in[0,M]).
\]
Choose
\[
\eta_{\mathrm{pert}}
=\min\left\{
\frac{m_{\mathrm{ext}}}{B_{\mathrm{sign}}^2},
\frac{\mathcal E-b_{\mathrm{ext}}}{2B_{\mathrm{sign}}}
\right\}>0.
\]
On the first region,
\[
\begin{aligned}
\{r_{\mathrm{best}}-\eta_{\mathrm{pert}}v_{\mathrm{sign}}\}^2
&\le
\mathcal E^2-2\eta_{\mathrm{pert}}m_{\mathrm{ext}}
+\eta_{\mathrm{pert}}^2B_{\mathrm{sign}}^2\\
&\le\mathcal E^2-\eta_{\mathrm{pert}}m_{\mathrm{ext}}
<\mathcal E^2.
\end{aligned}
\]
On its complement,
\[
|r_{\mathrm{best}}-\eta_{\mathrm{pert}}v_{\mathrm{sign}}|
<
b_{\mathrm{ext}}+\frac{\mathcal E-b_{\mathrm{ext}}}{2}
<\mathcal E.
\]
Continuity on the compact interval makes the supremum strictly smaller than \(\mathcal E\), contradicting optimality of
\(P_{\mathrm{best}}+\eta_{\mathrm{pert}}Q_{\mathrm{sign}}\). This proves the alternating-node assertion.
% lean: CausalSmith.Stat.OptvalueVanishingoverlapRate.weightedApproxError_alternating_witness

\item
For these nodes, define the Lagrange weights and their two normalizations by
\[
\lambda_\ell^{\mathrm{Lag}}
=\prod_{\substack{0\le r\le K+1\\r\ne\ell}}
(z_\ell^{\mathrm{dual}}-z_r^{\mathrm{dual}})^{-1},
\qquad
\lambda_\ell^{\mathrm{norm}}
=\frac{\lambda_\ell^{\mathrm{Lag}}}
       {\sum_{r=0}^{K+1}|\lambda_r^{\mathrm{Lag}}|},
\]
\[
D_{\mathrm{dual}}
=\sum_{\ell=0}^{K+1}
|\lambda_\ell^{\mathrm{norm}}|(1+z_\ell^{\mathrm{dual}}),
\qquad
\lambda_\ell^{\mathrm{dual}}
=\frac{\lambda_\ell^{\mathrm{norm}}}{D_{\mathrm{dual}}}.
\]
The node differences are nonzero, and the two normalization denominators are positive; moreover, \(D_{\mathrm{dual}}\ge1\).
Comparing the coefficient of degree \(K+1\) in the Lagrange interpolation formula gives
\[
\sum_{\ell=0}^{K+1}
\lambda_\ell^{\mathrm{norm}}P_{\mathrm{test}}(z_\ell^{\mathrm{dual}})
=0
\qquad(\deg P_{\mathrm{test}}\le K).
\]
The ordered nodes also give
\[
\lambda_\ell^{\mathrm{norm}}
=(-1)^{K+1-\ell}|\lambda_\ell^{\mathrm{norm}}|.
\]
Combining this sign with the alternating residual yields
\[
\begin{aligned}
\sum_{\ell=0}^{K+1}
\lambda_\ell^{\mathrm{norm}}\phi_\epsilon(z_\ell^{\mathrm{dual}})
&=
\sum_{\ell=0}^{K+1}
\lambda_\ell^{\mathrm{norm}}
\{\phi_\epsilon(z_\ell^{\mathrm{dual}})
-P_{\mathrm{best}}(z_\ell^{\mathrm{dual}})\}\\
&=o_{\mathrm{dual}}(-1)^{K+1}\mathcal E D_{\mathrm{dual}}.
\end{aligned}
\]
Consequently,
\[
\sum_{\ell=0}^{K+1}
\lambda_\ell^{\mathrm{dual}}(z_\ell^{\mathrm{dual}})^r=0
\quad(r=0,\ldots,K),
\]
\[
\sum_{\ell=0}^{K+1}
|\lambda_\ell^{\mathrm{dual}}|(1+z_\ell^{\mathrm{dual}})=1,
\qquad
\left|
\sum_{\ell=0}^{K+1}
\lambda_\ell^{\mathrm{dual}}\phi_\epsilon(z_\ell^{\mathrm{dual}})
\right|=\mathcal E.
\]
Split the weights into positive and negative parts:
\[
\sigma_+^{\mathrm{dual}}
=\sum_{\ell=0}^{K+1}(\lambda_\ell^{\mathrm{dual}})_+
 \delta_{z_\ell^{\mathrm{dual}}},
\qquad
\sigma_-^{\mathrm{dual}}
=\sum_{\ell=0}^{K+1}(-\lambda_\ell^{\mathrm{dual}})_+
 \delta_{z_\ell^{\mathrm{dual}}}.
\]
Define
\[
t_{\mathrm{dual}}=\sigma_+^{\mathrm{dual}}(\mathbb R),
\qquad
u_{\mathrm{dual}}=\int z\,d\sigma_+^{\mathrm{dual}},
\qquad
z_{\mathrm{dual}}^{\mathrm{comp}}
=\frac{1-u_{\mathrm{dual}}}{1-t_{\mathrm{dual}}}.
\]
The degree-zero and degree-one identities and the weighted normalization give
\[
t_{\mathrm{dual}}+u_{\mathrm{dual}}=\frac12,
\qquad
0\le t_{\mathrm{dual}}\le\frac12,
\qquad
\frac12\le z_{\mathrm{dual}}^{\mathrm{comp}}\le2.
\]
The same mass, mean, and moment calculation used for the packet completion shows that
\[
\nu_0^{\mathrm{dual}}
=\sigma_-^{\mathrm{dual}}
 +(1-t_{\mathrm{dual}})\delta_{z_{\mathrm{dual}}^{\mathrm{comp}}},
\qquad
\nu_1^{\mathrm{dual}}
=\sigma_+^{\mathrm{dual}}
 +(1-t_{\mathrm{dual}})\delta_{z_{\mathrm{dual}}^{\mathrm{comp}}}
\]
form a pair in \(\mathfrak N_{K,M}\), with target gap \(\mathcal E\). Hence
\[
\mathcal E\le\Delta.
\]
The feasible set is nonempty, since \((\delta_1,\delta_1)\) is feasible. Its gaps are bounded by \(4e(0)\), by the polynomial residual bound above. Taking the polynomial infimum in that bound, followed by the feasible-pair supremum, therefore gives
\[
\Delta\le4\mathcal E.
\]
% lean: CausalSmith.Stat.OptvalueVanishingoverlapRate.weightedApproxError_le_constrainedPriorSeparation

\item
For the approximation upper bound, the square-root coordinate gives the required uniform angular regularity. Define, for \(\xi\ge0\),
\[
b_{\mathrm{lo}}(\xi)=(1-\xi^2)_+,
\qquad
b_{\mathrm{hi}}(\xi)=\frac{(\xi^2-1)_+}{1+\xi^2}.
\]
For \(0\le\xi\le\eta\), splitting at \(1\) proves
\[
|b_{\mathrm{lo}}(\xi)-b_{\mathrm{lo}}(\eta)|
\le2(\eta-\xi),
\qquad
|b_{\mathrm{hi}}(\xi)-b_{\mathrm{hi}}(\eta)|
\le2(\eta-\xi).
\]
For the lower branch, when \(\eta\le1\), the difference equals
\((\eta-\xi)(\eta+\xi)\); when \(\xi\le1\le\eta\), it equals
\(1-\xi^2\le2(1-\xi)\); when \(1\le\xi\), it vanishes.
For the upper branch, when \(1\le\xi\),
\[
b_{\mathrm{hi}}(\eta)-b_{\mathrm{hi}}(\xi)
=\frac{2(\eta-\xi)(\eta+\xi)}
       {(1+\xi^2)(1+\eta^2)}
\le2(\eta-\xi).
\]
When \(\xi\le1\le\eta\), the bound follows from
\[
\frac{\eta^2-1}{1+\eta^2}\le2(\eta-1)\le2(\eta-\xi);
\]
when \(\eta\le1\), the difference vanishes. Interchanging \(\xi,\eta\) covers the opposite ordering. It follows that
\[
|g_\epsilon(\xi^2)-g_\epsilon(\eta^2)|
\le
2\{\epsilon+(1-2\epsilon)\}|\xi-\eta|
\le2|\xi-\eta|.
\]
Using the half-angle identity
\[
z_M(\omega)=\{\sqrt M\,|\cos(\omega/2)|\}^2
\]
and the Lipschitz bound for cosine gives
\[
|f_{M,\epsilon}(\omega)-f_{M,\epsilon}(\upsilon)|
\le\sqrt M\,|\omega-\upsilon|,
\qquad \omega,\upsilon\in\mathbb R.
\]

Set
\[
N_{\mathrm{up}}=\lfloor K/2\rfloor+1,
\qquad
T_{\mathrm{up}}(\omega)
=\int_{-\pi}^{\pi}
f_{M,\epsilon}(\omega+\upsilon)
J_{N_{\mathrm{up}}}^{\mathrm J}(\upsilon)\,d\upsilon.
\]
Convolution with the even Jackson kernel makes \(T_{\mathrm{up}}\) an even trigonometric polynomial of degree at most
\(2(N_{\mathrm{up}}-1)\le K\). Positivity, normalization, and the first kernel moment give
\[
|f_{M,\epsilon}(\omega)-T_{\mathrm{up}}(\omega)|
\le
\sqrt M\int_{-\pi}^{\pi}
|\upsilon|J_{N_{\mathrm{up}}}^{\mathrm J}(\upsilon)\,d\upsilon
\le\frac{32\sqrt M}{N_{\mathrm{up}}}.
\]
Writing its cosine modes as Chebyshev polynomials supplies a real polynomial \(Q_{\mathrm{up}}\) such that
\[
T_{\mathrm{up}}(\omega)=Q_{\mathrm{up}}(\cos\omega),
\qquad
\deg Q_{\mathrm{up}}\le2(N_{\mathrm{up}}-1).
\]
Define
\[
P_{\mathrm{up}}(z)
=\epsilon(z-1)+Q_{\mathrm{up}}(2z/M-1).
\]
Then \(\deg P_{\mathrm{up}}\le K\), and every \(z\in[0,M]\) can be written as
\(z_M(\omega)\). Since \(K\le2N_{\mathrm{up}}\),
\[
|\phi_\epsilon(z)-P_{\mathrm{up}}(z)|
\le\frac{32\sqrt M}{N_{\mathrm{up}}}
\le\frac{64\sqrt M}{K}.
\]
Division by \(1+z\ge1\), followed by the polynomial infimum, proves
\[
\mathcal E\le\frac{64\sqrt M}{K}.
\]
% lean: CausalSmith.Stat.OptvalueVanishingoverlapRate.weightedApproxError_upper

\item
Choose the universal constants
\[
c=\min\{\min(c_{\mathrm E},c_{\mathrm p}),1\},
\qquad
C=257.
\]
They satisfy \(0<c<C\), and the preceding bounds yield
\[
c\frac{\sqrt M}{K}
\le\mathcal E
\le\Delta
\le4\mathcal E
\le256\frac{\sqrt M}{K}
\le C\frac{\sqrt M}{K}.
\]
The packet completion retains its gap because \(c\le c_{\mathrm p}\).
All the packet estimates used the closed range
\(0\le\epsilon\le1/2\) and the non-strict support bound \(M\le K^2\), so they include both overlap endpoints and \(M=K^2\).
% lean: CausalSmith.Stat.OptvalueVanishingoverlapRate.weighted_approximation_and_packet

\item
\textbf{Sparse comparison.}\quad
We next establish an observed minimax lower bound for all alphabets. For this part, assume
\[
n\ge1,\qquad d\ge2,\qquad 0<\epsilon\le\tfrac12,
\]
and define
\[
L_d=\log(ed),
\qquad
r_{n,d,\epsilon}=\frac{d}{n\epsilon L_d},
\qquad
\rho_{n,d,\epsilon}=\min\{1,r_{n,d,\epsilon}\}.
\]
In particular \(L_d=1+\log d\ge1\) and \(\rho_{n,d,\epsilon}\ge0\).

Here is the likelihood calibration used in both large-alphabet constructions. For integers \(K_{\mathrm{test}}\ge1\) and real \(\tau\ge0\), define
\[
\mathcal T_{K_{\mathrm{test}}}(\tau)
=\sum_{\ell>K_{\mathrm{test}}}\frac{\tau^\ell}{\ell!},
\qquad
\beta_{\mathrm{tail}}=\frac1{16e}.
\]
When \(\tau\le\beta_{\mathrm{tail}}K_{\mathrm{test}}\), the factorial bound
\(\ell!\ge(\ell/e)^\ell\) gives, for \(\ell>K_{\mathrm{test}}\),
\[
\frac{\tau^\ell}{\ell!}
\le\left(\frac{e\tau}{\ell}\right)^\ell
\le4^{-\ell}.
\]
Summing the geometric tail therefore gives
\[
\mathcal T_{K_{\mathrm{test}}}(\tau)
\le4^{-K_{\mathrm{test}}},
\qquad
\sqrt{\mathcal T_{K_{\mathrm{test}}}(\tau)}
\le2^{-K_{\mathrm{test}}}.
\]
Define
\[
C_{\mathrm{tv}}=\frac{|\log8|+2}{\log2}.
\]
If \(K_{\mathrm{test}}\ge C_{\mathrm{tv}}L_d\), then
\[
K_{\mathrm{test}}\log2
\ge(|\log8|+2)(1+\log d)
\ge\log(8d),
\]
and consequently
\[
d\sqrt{\mathcal T_{K_{\mathrm{test}}}(\tau)}
\le d\,2^{-K_{\mathrm{test}}}\le\frac18.
\]
% lean: CausalSmith.Stat.OptvalueVanishingoverlapRate.sparse_geometric_degree_calibration

\item
Consider first the sparse experiment of
\cref{obj:lem:sparse-likelihood}, with inflation \(B>2\) and a constrained pair of degree \(K_{\mathrm{test}}\). Under prior \(\nu_\iota\), \(\iota\in\{0,1\}\), draw cell intensities independently with that prior, and observe all four independent Poisson counts in each cell.

For one cell, let \(\mathbb Q_\star\) be the reference count law at \(z_\star=M/2\), and define the mixture densities
\[
h_{\mathrm{sp},\iota}
=\int L_z\,d\nu_\iota(z),
\qquad \iota\in\{0,1\}.
\]
The likelihood ratios \(L_z\) and Gram coefficient \(\alpha\) are those in
\cref{obj:lem:sparse-likelihood}; at a zero alternative mean, the count factor uses \(0^0=1\). Its Gram identity, followed by expansion of the exponential, gives
\[
\begin{aligned}
\int(h_{\mathrm{sp},1}-h_{\mathrm{sp},0})^2\,d\mathbb Q_\star
&=
\sum_{\ell>K_{\mathrm{test}}}\frac{\alpha^\ell}{\ell!}
\left\{
\int(z-z_\star)^\ell\,d(\nu_1-\nu_0)(z)
\right\}^2\\
&\le
4\mathcal T_{K_{\mathrm{test}}}(\tau_{\mathrm{sp}}),
\qquad
\tau_{\mathrm{sp}}=\frac{\alpha M^2}{4}.
\end{aligned}
\]
The terms through degree \(K_{\mathrm{test}}\) vanish by moment matching, and the remaining moments have absolute value at most
\(2(M/2)^\ell\). Thus Cauchy--Schwarz and telescoping products imply
\[
\operatorname{TV}
\left(\mathbb Q_{\mathrm{sp},0}^{\otimes d},
      \mathbb Q_{\mathrm{sp},1}^{\otimes d}\right)
\le
d\sqrt{\mathcal T_{K_{\mathrm{test}}}(\tau_{\mathrm{sp}})},
\]
where
\[
\mathbb Q_{\mathrm{sp},\iota}
=h_{\mathrm{sp},\iota}\,\mathbb Q_\star.
\]
Set
\[
\eta_{\mathrm{sp}}=2\beta_{\mathrm{tail}}.
\]
The bound on \(\alpha\) in
\cref{obj:lem:sparse-likelihood} gives
\[
\tau_{\mathrm{sp}}\le\frac{Bn\epsilon M}{2d}.
\]
Hence
\[
M\le\frac{\eta_{\mathrm{sp}}dK_{\mathrm{test}}}{Bn\epsilon},
\qquad
K_{\mathrm{test}}\ge C_{\mathrm{tv}}L_d
\]
ensure that the full product-mixture total variation is at most \(1/8\).
% lean: CausalSmith.Stat.OptvalueVanishingoverlapRate.sparse_product_mixture_tv_le_one_eighth

\item
For either sparse product prior, write the independent cell intensities as
\[
\mathbf Z^{\mathrm{sp}}
=(Z_x^{\mathrm{sp}})_{x\in[d]}
\sim\nu_\iota^{\otimes d},
\qquad
\mathbf Z^{\mathrm{sp}}\in[0,M]^d,
\]
and define
\[
H_{\mathrm{sp}}
=\frac1d\sum_{x\in[d]}\phi_\epsilon(Z_x^{\mathrm{sp}}),
\qquad
S_{\mathrm{sp}}
=1-\epsilon+\frac{\epsilon}{d}\sum_{x\in[d]}Z_x^{\mathrm{sp}},
\]
\[
\theta_{\mathrm{sp},\iota}
=\frac12+\frac12\int\phi_\epsilon\,d\nu_\iota,
\qquad
\delta_{\mathrm{sp}}
=|\theta_{\mathrm{sp},1}-\theta_{\mathrm{sp},0}|
=\frac12\left|\int\phi_\epsilon\,d(\nu_1-\nu_0)\right|.
\]
The sparse value identity in
\cref{obj:lem:sparse-likelihood} is
\[
\Psi(\mathbb P_{\mathbf Z^{\mathrm{sp}}}^{\mathrm{sp}})
=\frac12+\frac{H_{\mathrm{sp}}}{2S_{\mathrm{sp}}}.
\]
The mean-one constraint and \(0\le Z_x^{\mathrm{sp}}\le M\) give
\[
E (Z_x^{\mathrm{sp}})^2\le M,
\qquad
E S_{\mathrm{sp}}=1,
\qquad
E(S_{\mathrm{sp}}-1)^2\le\frac{\epsilon^2M}{d}.
\]
For \(0\le z\le M\), both \(z-1\) and zero are at most \(1+z\), so
\[
0\le(z-1)_+\le1+z,
\qquad
0\le1-\epsilon+\epsilon z,
\qquad
1+z>0.
\]
Multiplying the first inequality by the nonnegative factor
\((1-\epsilon+\epsilon z)/(1+z)\) gives
\[
0\le\phi_\epsilon(z)\le1-\epsilon+\epsilon z.
\]
Moreover, \(z^2\le Mz\), \((1-\epsilon)^2\le1\), and the nonnegativity of
\((1-\epsilon-\epsilon z)^2\) yields
\[
\begin{aligned}
\phi_\epsilon(z)^2
&\le(1-\epsilon+\epsilon z)^2\\
&\le2\{(1-\epsilon)^2+\epsilon^2z^2\}\\
&\le2(1+\epsilon^2Mz).
\end{aligned}
\]
Independence therefore yields
\[
E\{H_{\mathrm{sp}}-EH_{\mathrm{sp}}\}^2
\le\frac{2(1+\epsilon^2M)}d.
\]
Since \(S_{\mathrm{sp}}\ge1-\epsilon\ge1/2\) and
\(0\le H_{\mathrm{sp}}\le S_{\mathrm{sp}}\),
\[
\left|\frac{H_{\mathrm{sp}}}{S_{\mathrm{sp}}}
-H_{\mathrm{sp}}\right|
\le|1-S_{\mathrm{sp}}|.
\]
Using \((a+b)^2\le2(a^2+b^2)\) now gives
\[
E\left[
\left\{\Psi(\mathbb P_{\mathbf Z^{\mathrm{sp}}}^{\mathrm{sp}})
-\theta_{\mathrm{sp},\iota}\right\}^2
\right]
\le\frac{2+3\epsilon^2M}{2d}.
\]
Consequently, if
\[
64(2+3\epsilon^2M)\le d\delta_{\mathrm{sp}}^2,
\]
Markov's inequality gives
\[
\Pr_\iota\left(
\left|\Psi(\mathbb P_{\mathbf Z^{\mathrm{sp}}}^{\mathrm{sp}})
-\theta_{\mathrm{sp},\iota}\right|
>\frac{\delta_{\mathrm{sp}}}{4}
\right)\le\frac18.
\]

For completeness, define the two integrated squared risks of a count estimator \(\widehat W\) by
\[
\mathcal B_{\mathrm{sp},\iota}(\widehat W)
=
E_{\mathbf Z^{\mathrm{sp}}\sim\nu_\iota^{\otimes d}}
E_{\mathrm{counts}\mid\mathbf Z^{\mathrm{sp}}}
\left[
\{\widehat W-\Psi(\mathbb P_{\mathbf Z^{\mathrm{sp}}}^{\mathrm{sp}})\}^2
\right].
\]
If \(\theta_{\mathrm{sp},0}\le\theta_{\mathrm{sp},1}\), threshold
\(\widehat W\) at
\[
\frac{\theta_{\mathrm{sp},0}+\theta_{\mathrm{sp},1}}2.
\]
Otherwise interchange the prior indices and use the same midpoint. On an erroneous decision with the target within
\(\delta_{\mathrm{sp}}/4\) of its center, estimation error is at least
\(\delta_{\mathrm{sp}}/4\). The sum of the two testing errors is at least one minus their predictive total variation. Subtracting the two target-deviation probabilities thus proves
\[
\begin{aligned}
\max_{\iota=0,1}\mathcal B_{\mathrm{sp},\iota}(\widehat W)
&\ge
\frac12\left(\frac{\delta_{\mathrm{sp}}}{4}\right)^2
\left(1-\frac18-\frac18-\frac18\right)\\
&=\frac{5\delta_{\mathrm{sp}}^2}{256}.
\end{aligned}
\]

Use the completed packet pair, whose gap is at least
\(c_{\mathrm p}\sqrt M/K_{\mathrm{test}}\). Then
\[
\delta_{\mathrm{sp}}^2
\ge\frac{c_{\mathrm p}^2M}{4K_{\mathrm{test}}^2}.
\]
Thus the condition
\[
256(2+3\epsilon^2M)K_{\mathrm{test}}^2
\le c_{\mathrm p}^2Md
\]
supplies the target-deviation bound, and, with
\[
c_{\mathrm{sp}}=\frac{5c_{\mathrm p}^2}{1024},
\]
every measurable count estimator satisfies
\[
\max_{\iota=0,1}\mathcal B_{\mathrm{sp},\iota}(\widehat W)
\ge c_{\mathrm{sp}}\frac{M}{K_{\mathrm{test}}^2}.
\]
% lean: CausalSmith.Stat.OptvalueVanishingoverlapRate.sparse_packet_bayesRisk_lower

\item
We record the logarithmic degree and alphabet calibration explicitly. For real \(C_{\mathrm{deg}}\ge2\) and integer \(d\ge2\), define
\[
\mathcal K_{\mathrm{cal}}(C_{\mathrm{deg}},d)
=2\left\lceil\frac{C_{\mathrm{deg}}L_d}{2}\right\rceil.
\]
The ceiling inequalities and \(L_d\ge1\) give
\[
2\le\mathcal K_{\mathrm{cal}}(C_{\mathrm{deg}},d),
\qquad
C_{\mathrm{deg}}L_d
\le\mathcal K_{\mathrm{cal}}(C_{\mathrm{deg}},d)
\le(C_{\mathrm{deg}}+2)L_d.
\]
Choose
\[
C_{\mathrm{sp}}=\max\{C_{\mathrm{tv}},2\},
\qquad
K_{\mathrm{sp}}=\mathcal K_{\mathrm{cal}}(C_{\mathrm{sp}},d).
\]
Because \(L_d^2/d\to0\), a universal integer \(D_{\mathrm{sp}}\ge2\) can be chosen so that, for \(d\ge D_{\mathrm{sp}}\),
\[
\frac{L_d^2}{d}
<
\frac{c_{\mathrm p}^2}{448(C_{\mathrm{sp}}+2)^2},
\qquad
448K_{\mathrm{sp}}^2\le c_{\mathrm p}^2d.
\]
The asserted logarithmic limit follows, for example, by writing
\(d=e^{L_d-1}\) and using that an exponential dominates a fixed power. For \(M\ge2\) and \(0\le\epsilon\le1/2\),
\[
256(2+3\epsilon^2M)\le448M.
\]
Hence this single cutoff ensures
\[
256(2+3\epsilon^2M)K_{\mathrm{sp}}^2
\le c_{\mathrm p}^2Md
\]
uniformly in the required support and overlap ranges. The previous two steps therefore apply whenever
\[
2\le M\le K_{\mathrm{sp}}^2,
\qquad
M\le\frac{\eta_{\mathrm{sp}}dK_{\mathrm{sp}}}{Bn\epsilon}.
\]
% lean: CausalSmith.Stat.OptvalueVanishingoverlapRate.sparse_packet_bayesRisk_lower_largeAlphabet

\item
\textbf{Dense comparison.}\quad
The dense construction uses symmetric moment-matched priors on \([-1,1]\). We establish the absolute-moment gap used for them.

For \(K_{\mathrm{abs}}\ge1\), define
\[
E_{\mathrm{abs}}(K_{\mathrm{abs}})
=\inf_{\deg P_{\mathrm{abs}}\le K_{\mathrm{abs}}}
\sup_{-1\le z\le1}\bigl||z|-P_{\mathrm{abs}}(z)\bigr|.
\]
On the circle of period \(\pi\), define, for each integer \(N_{\mathrm{Fej}}\ge1\),
\[
F_{\mathrm{Fej},N_{\mathrm{Fej}}}(\omega)
=\frac1{N_{\mathrm{Fej}}}
\left|\sum_{\ell_{\mathrm{Fej}}=0}^{N_{\mathrm{Fej}}-1}
 e^{2i\ell_{\mathrm{Fej}}\omega}\right|^2,
\qquad \omega\in\mathbb R,
\]
\[
V_{\mathrm{VP},K_{\mathrm{abs}}}
=2F_{\mathrm{Fej},2K_{\mathrm{abs}}}
-F_{\mathrm{Fej},K_{\mathrm{abs}}}.
\]
For a continuous complex-valued \(\pi\)-periodic function \(f_{\mathrm{cusp}}\), define
\[
\Lambda_{\mathrm{cusp},K_{\mathrm{abs}}}(f_{\mathrm{cusp}})
=f_{\mathrm{cusp}}(0)-\int_0^\pi
V_{\mathrm{VP},K_{\mathrm{abs}}}(\omega)
 f_{\mathrm{cusp}}(\omega)\,\frac{d\omega}{\pi}.
\]
Expanding the squared finite sum gives
\[
F_{\mathrm{Fej},N_{\mathrm{Fej}}}(\omega)
=\frac1{N_{\mathrm{Fej}}}
\sum_{\ell_{\mathrm{Fej}},r_{\mathrm{Fej}}=0}^{N_{\mathrm{Fej}}-1}
 e^{2i(\ell_{\mathrm{Fej}}-r_{\mathrm{Fej}})\omega}.
\]
For each integer frequency \(j_{\mathrm{Four}}\), orthogonality retains exactly the pairs whose index difference is \(-j_{\mathrm{Four}}\). There are
\(N_{\mathrm{Fej}}-|j_{\mathrm{Four}}|\) such pairs when
\(|j_{\mathrm{Four}}|<N_{\mathrm{Fej}}\), and none otherwise. Thus the triangular coefficients are
\[
\begin{aligned}
\gamma_{\mathrm{Fej}}(N_{\mathrm{Fej}},j_{\mathrm{Four}})
&:=\int_0^\pi F_{\mathrm{Fej},N_{\mathrm{Fej}}}(\omega)
 e^{2ij_{\mathrm{Four}}\omega}\,\frac{d\omega}{\pi}\\
&=\begin{cases}
1-|j_{\mathrm{Four}}|/N_{\mathrm{Fej}},
&|j_{\mathrm{Four}}|<N_{\mathrm{Fej}},\\
0,&|j_{\mathrm{Four}}|\ge N_{\mathrm{Fej}}.
\end{cases}
\end{aligned}
\]
In particular, both Fej\'er kernels are nonnegative with integral one, so
\[
|\Lambda_{\mathrm{cusp},K_{\mathrm{abs}}}(f_{\mathrm{cusp}})|
\le4\|f_{\mathrm{cusp}}\|_\infty.
\]
Substituting the triangular coefficients into the cusp functional gives
\[
\begin{aligned}
\mu_{\mathrm{cusp}}(j_{\mathrm{Four}})
&:=\Lambda_{\mathrm{cusp},K_{\mathrm{abs}}}
 (\omega\mapsto e^{2ij_{\mathrm{Four}}\omega})\\
&=1-2\gamma_{\mathrm{Fej}}(2K_{\mathrm{abs}},j_{\mathrm{Four}})
 +\gamma_{\mathrm{Fej}}(K_{\mathrm{abs}},j_{\mathrm{Four}})\\
&=\begin{cases}
0,&|j_{\mathrm{Four}}|\le K_{\mathrm{abs}},\\
|j_{\mathrm{Four}}|/K_{\mathrm{abs}}-1,
&K_{\mathrm{abs}}<|j_{\mathrm{Four}}|<2K_{\mathrm{abs}},\\
1,&|j_{\mathrm{Four}}|\ge2K_{\mathrm{abs}}.
\end{cases}
\end{aligned}
\]
These multipliers are nonnegative. The identity
\[
\sin2\omega=\frac{e^{2i\omega}-e^{-2i\omega}}{2i}
\]
and the binomial expansion show that a polynomial of degree at most
\(K_{\mathrm{abs}}\) in \(\sin2\omega\) has frequencies
\(|j_{\mathrm{Four}}|\le K_{\mathrm{abs}}\). Hence
\[
\Lambda_{\mathrm{cusp},K_{\mathrm{abs}}}
 (\omega\mapsto P_{\mathrm{abs}}(\sin2\omega))=0
\qquad(\deg P_{\mathrm{abs}}\le K_{\mathrm{abs}}).
\]

To compute the absolute-sine coefficients, define
\[
S_{\mathrm{abs}}(\omega)=|\sin\omega|,
\qquad
c_{\mathrm{Four},\mathrm{abs}}(j_{\mathrm{Four}})
=\int_0^\pi S_{\mathrm{abs}}(\omega)e^{-2ij_{\mathrm{Four}}\omega}
 \,\frac{d\omega}{\pi},
\qquad j_{\mathrm{Four}}\in\mathbb Z.
\]
On \([0,\pi]\), \(S_{\mathrm{abs}}(\omega)=\sin\omega\). For every integer
\(j_{\mathrm{Four}}\), the denominator \(1-4j_{\mathrm{Four}}^2\) is nonzero, and the function
\[
\mathcal A_{\mathrm{Four},j_{\mathrm{Four}}}(\omega)
=e^{-2ij_{\mathrm{Four}}\omega}
\frac{-2ij_{\mathrm{Four}}\sin\omega-\cos\omega}
 {1-4j_{\mathrm{Four}}^2},
\qquad \omega\in[0,\pi],
\]
satisfies
\[
\mathcal A_{\mathrm{Four},j_{\mathrm{Four}}}'(\omega)
=e^{-2ij_{\mathrm{Four}}\omega}\sin\omega.
\]
Evaluating at the endpoints, using
\(e^{-2ij_{\mathrm{Four}}\pi}=1\), gives
\[
c_{\mathrm{Four},\mathrm{abs}}(j_{\mathrm{Four}})
=\frac{\mathcal A_{\mathrm{Four},j_{\mathrm{Four}}}(\pi)
 -\mathcal A_{\mathrm{Four},j_{\mathrm{Four}}}(0)}{\pi}
=-\frac2{\pi(4j_{\mathrm{Four}}^2-1)}.
\]
The zero coefficient is \(2/\pi\). For \(|j_{\mathrm{Four}}|\ge1\),
\[
|c_{\mathrm{Four},\mathrm{abs}}(j_{\mathrm{Four}})|
=\frac2{\pi(4j_{\mathrm{Four}}^2-1)}
\le\frac1{j_{\mathrm{Four}}^2},
\qquad
\sum_{j_{\mathrm{Four}}\in\mathbb Z}
 |c_{\mathrm{Four},\mathrm{abs}}(j_{\mathrm{Four}})|<\infty.
\]
Thus the Fourier series converges absolutely and uniformly; its coefficients identify its sum with the continuous function \(S_{\mathrm{abs}}\). Pairing opposite frequencies and replacing \(\omega\) by \(2\omega\) yields
\[
|\sin2\omega|
=\frac2\pi-\frac4\pi
 \sum_{\ell=1}^\infty\frac{\cos(4\ell\omega)}{4\ell^2-1}.
\]
Continuity of the cusp functional permits termwise application, giving
\[
-\Lambda_{\mathrm{cusp},K_{\mathrm{abs}}}(|\sin2\cdot|)
=\sum_{\ell=1}^\infty
 \frac{4\mu_{\mathrm{cusp}}(2\ell)}{\pi(4\ell^2-1)}.
\]
Every summand is nonnegative. For
\(K_{\mathrm{abs}}\le\ell<2K_{\mathrm{abs}}\), the multiplier equals one, and \(\pi\le4\) gives
\[
\frac4{\pi(4\ell^2-1)}\ge\frac1{16K_{\mathrm{abs}}^2}.
\] Keeping the
\(K_{\mathrm{abs}}\) terms with
\(K_{\mathrm{abs}}\le\ell<2K_{\mathrm{abs}}\) gives
\[
-\Lambda_{\mathrm{cusp},K_{\mathrm{abs}}}(|\sin2\cdot|)
\ge
\sum_{\ell=K_{\mathrm{abs}}}^{2K_{\mathrm{abs}}-1}
\frac4{\pi(4\ell^2-1)}
\ge\frac1{16K_{\mathrm{abs}}}.
\]
It follows that
\[
\frac1{16K_{\mathrm{abs}}}
\le
4\sup_{-1\le z\le1}\bigl||z|-P_{\mathrm{abs}}(z)\bigr|,
\]
and therefore
\[
E_{\mathrm{abs}}(K_{\mathrm{abs}})
\ge\frac1{100K_{\mathrm{abs}}}.
\]

To construct the priors, let
\[
\mathcal V_{\mathrm{abs}}
=\{P_{\mathrm{abs}}|_{[-1,1]}:\deg P_{\mathrm{abs}}\le K_{\mathrm{abs}}\}.
\]
The norm of the class of \(|z|\) in
\(C([-1,1],\mathbb R)/\mathcal V_{\mathrm{abs}}\) is
\(E_{\mathrm{abs}}(K_{\mathrm{abs}})\). A norming functional on that quotient, pulled back to \(C([-1,1],\mathbb R)\), supplies a functional
\(\Lambda_{\mathrm{abs}}\) satisfying
\[
\|\Lambda_{\mathrm{abs}}\|\le1,
\qquad
\Lambda_{\mathrm{abs}}(z^\ell)=0
\quad(0\le\ell\le K_{\mathrm{abs}}),
\qquad
\Lambda_{\mathrm{abs}}(|z|)=E_{\mathrm{abs}}(K_{\mathrm{abs}}).
\]
Its positive half-split can be obtained by positive extension. In detail, let \(1\) denote the constant function with value one and consider the cone
\[
\mathcal C_{\mathrm{split}}
=\{(f_++f_-,-\Lambda_{\mathrm{abs}}(f_-)):
 f_+,f_-\in C([-1,1],\mathbb R),\ f_+,f_-\ge0\}.
\]
On the subspace of pairs \((a_{\mathrm{cone}}\,1,b_{\mathrm{cone}})\), prescribe the functional
\[
(a_{\mathrm{cone}}\,1,b_{\mathrm{cone}})
\longmapsto a_{\mathrm{cone}}/2+b_{\mathrm{cone}},
\qquad a_{\mathrm{cone}},b_{\mathrm{cone}}\in\mathbb R.
\]
If such a pair belongs to the cone, then
\[
f_++f_-=a_{\mathrm{cone}}\,1,
\qquad
b_{\mathrm{cone}}=-\Lambda_{\mathrm{abs}}(f_-),
\qquad a_{\mathrm{cone}}\ge0.
\]
Consequently,
\[
\|2f_--a_{\mathrm{cone}}\,1\|_\infty\le a_{\mathrm{cone}},
\qquad
\Lambda_{\mathrm{abs}}(f_-)
=\tfrac12\Lambda_{\mathrm{abs}}(2f_--a_{\mathrm{cone}}\,1)
\le a_{\mathrm{cone}}/2,
\]
where the last bound uses \(\|\Lambda_{\mathrm{abs}}\|\le1\) and
\(\Lambda_{\mathrm{abs}}(1)=0\). Thus the prescribed functional is nonnegative on its intersection with the cone.

The subspace also satisfies the required majorization condition. For every
\(f_{\mathrm{split}}\in C([-1,1],\mathbb R)\) and
\(t_{\mathrm{split}}\in\mathbb R\),
\[
(\|f_{\mathrm{split}}\|_\infty\,1,-t_{\mathrm{split}})
 +(f_{\mathrm{split}},t_{\mathrm{split}})
=(\|f_{\mathrm{split}}\|_\infty\,1+f_{\mathrm{split}},0)
\in\mathcal C_{\mathrm{split}}.
\]
Indeed, \(\|f_{\mathrm{split}}\|_\infty\,1+f_{\mathrm{split}}\ge0\), so this is the cone element obtained with that function as \(f_+\) and zero as \(f_-\). Positive extension therefore supplies a linear functional
\[
G_{\mathrm{split}}:C([-1,1],\mathbb R)\times\mathbb R\longrightarrow\mathbb R
\]
that is nonnegative on \(\mathcal C_{\mathrm{split}}\) and satisfies
\[
G_{\mathrm{split}}(a_{\mathrm{cone}}\,1,b_{\mathrm{cone}})
=a_{\mathrm{cone}}/2+b_{\mathrm{cone}}.
\]
Define, for every \(f_{\mathrm{split}}\in C([-1,1],\mathbb R)\),
\[
\Lambda_+(f_{\mathrm{split}})=G_{\mathrm{split}}(f_{\mathrm{split}},0),
\qquad
\Lambda_-(f_{\mathrm{split}})
=\Lambda_+(f_{\mathrm{split}})-\Lambda_{\mathrm{abs}}(f_{\mathrm{split}}).
\]
For \(f_{\mathrm{split}}\ge0\), both
\((f_{\mathrm{split}},0)\) and
\((f_{\mathrm{split}},-\Lambda_{\mathrm{abs}}(f_{\mathrm{split}}))\) belong to the cone. Since
\(G_{\mathrm{split}}(0,t_{\mathrm{split}})=t_{\mathrm{split}}\), this gives
\[
\Lambda_+(f_{\mathrm{split}})\ge0,
\qquad
\Lambda_-(f_{\mathrm{split}})
=G_{\mathrm{split}}(f_{\mathrm{split}},-\Lambda_{\mathrm{abs}}(f_{\mathrm{split}}))\ge0.
\]
Their masses and difference are
\[
\Lambda_+(1)=G_{\mathrm{split}}(1,0)=\frac12,
\qquad
\Lambda_-(1)=\Lambda_+(1)-\Lambda_{\mathrm{abs}}(1)=\frac12,
\qquad
\Lambda_+-\Lambda_-=\Lambda_{\mathrm{abs}}.
\]
Positivity bounds each functional on a continuous function by its mass times the uniform norm, so both have representing finite measures.
Their representing measures have mass \(1/2\). Averaging each measure with its reflection under \(z\mapsto-z\) preserves the absolute-value pairing and the annihilation of all powers through degree \(K_{\mathrm{abs}}\). Doubling these symmetric measures gives symmetric probability measures
\(\eta_0,\eta_1\) on \([-1,1]\) satisfying
\[
\int z^\ell\,d\eta_0=\int z^\ell\,d\eta_1
\quad(0\le\ell\le K_{\mathrm{abs}}),
\]
\[
\int|z|\,d\eta_1-\int|z|\,d\eta_0
=2E_{\mathrm{abs}}(K_{\mathrm{abs}})
\ge\frac1{50K_{\mathrm{abs}}}.
\]
% lean: CausalSmith.Stat.OptvalueVanishingoverlapRate.dense_symmetric_priors_absGap_lower

\item
For \(0<\mathfrak a\le1/2\), define the independent prior draws by
\[
\mathbf T^{\mathrm{de}}
=(T_x^{\mathrm{de}})_{x\in[d]}
\sim\eta_\iota^{\otimes d},
\qquad
\mathbf T^{\mathrm{de}}\in[-1,1]^d,
\]
and define the dense observed law by
\[
p_x=\frac1d,
\qquad
\pi_x=\epsilon,
\qquad
\mu_{0x}=\frac12,
\qquad
\mu_{1x}=\frac12+\mathfrak a T_x^{\mathrm{de}}.
\]
Equivalently, its four atom probabilities are
\[
q_{00,x}=q_{01,x}=\frac{1-\epsilon}{2d},
\qquad
q_{10,x}=\frac{\epsilon}{d}\left(\frac12-\mathfrak a T_x^{\mathrm{de}}\right),
\qquad
q_{11,x}=\frac{\epsilon}{d}\left(\frac12+\mathfrak a T_x^{\mathrm{de}}\right).
\]
These are legal observed laws, including the endpoints \(T_x^{\mathrm{de}}=\pm1\).
Their target and prior centers are
\[
\Psi_{\mathrm{de}}(\mathbf T^{\mathrm{de}})
=\frac12+\frac{\mathfrak a}{2d}
 \sum_{x\in[d]}(|T_x^{\mathrm{de}}|+T_x^{\mathrm{de}}),
\]
\[
\theta_{\mathrm{de},\iota}
=\frac12+\frac{\mathfrak a}{2}\int|z|\,d\eta_\iota,
\qquad
\delta_{\mathrm{de}}
=\frac{\mathfrak a}{2}
 \left(\int|z|\,d\eta_1-\int|z|\,d\eta_0\right)
\ge\frac{\mathfrak a}{100K_{\mathrm{abs}}}.
\]
Symmetry gives \(\int z\,d\eta_\iota=0\). Since
\((|z|+z)/2=z_+\in[0,1]\), independence gives
\[
E_\iota\{\Psi_{\mathrm{de}}-\theta_{\mathrm{de},\iota}\}^2
\le\frac{\mathfrak a^2}{d}.
\]

In the Poisson experiment, define the total treated intensity per cell by
\[
\lambda_{\mathrm{tr}}=\frac{Bn\epsilon}{d}.
\]
The treated counts have means
\[
\lambda_{\mathrm{tr}}\left(\frac12-\mathfrak a z\right),
\qquad
\lambda_{\mathrm{tr}}\left(\frac12+\mathfrak a z\right).
\]
For nonnegative integer treated failure and success counts, respectively denoted by \(N_{\mathrm{de},0}\) and \(N_{\mathrm{de},1}\), define the likelihood ratio relative to \(z=0\) by
\[
\begin{aligned}
L_{\mathrm{de},z}(N_{\mathrm{de},0},N_{\mathrm{de},1})
&=e^{\lambda_{\mathrm{tr}}\mathfrak a z}
 (1-2\mathfrak a z)^{N_{\mathrm{de},0}}
 e^{-\lambda_{\mathrm{tr}}\mathfrak a z}
 (1+2\mathfrak a z)^{N_{\mathrm{de},1}}\\
&=(1-2\mathfrak a z)^{N_{\mathrm{de},0}}
  (1+2\mathfrak a z)^{N_{\mathrm{de},1}},
\qquad z\in[-1,1].
\end{aligned}
\]
At a zero alternative mean, use \(0^0=1\); a zero base raised to a positive count is zero. The exponential factors cancel because the sum of the two treated means is \(\lambda_{\mathrm{tr}}\). Multiplication of these likelihood factors and their expectations gives the Gram identity
\[
E_0[L_{\mathrm{de},z}L_{\mathrm{de},w}]
=\exp\{4\lambda_{\mathrm{tr}}\mathfrak a^2zw\},
\qquad z,w\in[-1,1].
\]
Thus, defining
\[
\tau_{\mathrm{de}}=4\lambda_{\mathrm{tr}}\mathfrak a^2,
\qquad
\eta_{\mathrm{de}}=\frac{\beta_{\mathrm{tail}}}{4},
\]
the same moment-matching and Cauchy--Schwarz argument gives predictive total variation at most
\[
d\sqrt{\mathcal T_{K_{\mathrm{abs}}}(\tau_{\mathrm{de}})}
\le\frac18
\]
whenever
\[
K_{\mathrm{abs}}\ge C_{\mathrm{tv}}L_d,
\qquad
\lambda_{\mathrm{tr}}\mathfrak a^2
\le\eta_{\mathrm{de}}K_{\mathrm{abs}}.
\]
The two untreated counts have common means
\(Bn(1-\epsilon)/(2d)\), independent of the prior parameter. Adding them preserves this total-variation bound, so it applies to all four observed counts.

If
\[
128\mathfrak a^2\le d\delta_{\mathrm{de}}^2,
\]
the target-deviation probability at radius
\(\delta_{\mathrm{de}}/4\) is at most \(1/8\). Applying the midpoint argument above to the two dense centers gives, for every measurable estimator of the full counts,
\[
\max_{\iota=0,1}\mathcal B_{\mathrm{de},\iota}(\widehat W)
\ge\frac{5\delta_{\mathrm{de}}^2}{256}
\ge\frac5{2560000}\frac{\mathfrak a^2}{K_{\mathrm{abs}}^2},
\]
where
\[
\mathcal B_{\mathrm{de},\iota}(\widehat W)
=
E_{\mathbf T^{\mathrm{de}}\sim\eta_\iota^{\otimes d}}
E_{\mathrm{counts}\mid\mathbf T^{\mathrm{de}}}
[\{\widehat W-\Psi_{\mathrm{de}}(\mathbf T^{\mathrm{de}})\}^2].
\]
Choose
\[
C_{\mathrm{de}}=\max\{C_{\mathrm{tv}},2\},
\qquad
K_{\mathrm{de}}=\mathcal K_{\mathrm{cal}}(C_{\mathrm{de}},d).
\]
The same logarithmic cutoff construction, applied with packet comparison constant \(1/100\), supplies a universal
\(D_{\mathrm{de}}\ge2\) such that
\[
1280000K_{\mathrm{de}}^2\le d
\quad(d\ge D_{\mathrm{de}}).
\]
Since
\(\delta_{\mathrm{de}}^2\ge\mathfrak a^2/(10000K_{\mathrm{de}}^2)\),
this implies the required target-deviation condition. Reordering the four counts into the observed atom histogram is a bijection, so both integrated-risk bounds hold on that histogram space.
% lean: CausalSmith.Stat.OptvalueVanishingoverlapRate.dense_full_observed_logDegree_bayesRisk_lower

\item
We transfer these bounds to fixed samples. Let \(\widehat V\) be any measurable estimator based on \(n\) observations. For a histogram
\[
\boldsymbol h_{\mathrm{hist}}:\mathcal J_d\to\mathbb N,
\qquad
n_{\mathrm{hist}}=\sum_{o\in\mathcal J_d}\boldsymbol h_{\mathrm{hist}}(o),
\]
define
\[
\mathcal S(\boldsymbol h_{\mathrm{hist}})
=\{\boldsymbol s\in\mathcal J_d^{n_{\mathrm{hist}}}:
\boldsymbol s\text{ has histogram }\boldsymbol h_{\mathrm{hist}}\},
\qquad
\boldsymbol o^\circ=((0,0,0),\ldots,(0,0,0))\in\mathcal J_d^n,
\]
\[
\boldsymbol o_n(\boldsymbol s)=
\begin{cases}
(s_1,\ldots,s_n),&n_{\mathrm{hist}}\ge n,\\
\boldsymbol o^\circ,&n_{\mathrm{hist}}<n,
\end{cases}
\]
\[
\widehat W_{\widehat V}(\boldsymbol h_{\mathrm{hist}})
=\frac1{|\mathcal S(\boldsymbol h_{\mathrm{hist}})|}
 \sum_{\boldsymbol s\in\mathcal S(\boldsymbol h_{\mathrm{hist}})}
 \Pi_{[0,1]}\{\widehat V(\boldsymbol o_n(\boldsymbol s))\}.
\]
The set \(\mathcal S(\boldsymbol h_{\mathrm{hist}})\) is nonempty, including the zero histogram, for which it contains the empty string.

For independent Poisson atom counts with total mean
\(\lambda_{\mathrm{tot}}\), conditional on their total, the labels are independent with the normalized observed law; conditional on the histogram, their ordering is uniform on \(\mathcal S(h)\). Jensen's inequality, followed by projection onto \([0,1]\), therefore gives
\[
E\{\widehat W_{\widehat V}-\Psi\}^2
\le E\{\widehat V-\Psi\}^2
+\Pr\{\operatorname{Pois}(\lambda_{\mathrm{tot}})<n\}.
\]
For the sparse experiment,
\[
\lambda_{\mathrm{tot}}=BnS_{\mathrm{sp}}\ge Bn/2;
\]
for the dense experiment, \(\lambda_{\mathrm{tot}}=Bn\).

For any \(\lambda_{\mathrm{tot}}\ge Bn/2>n\), put
\[
\gamma_{\mathrm{sh}}
=\frac{\lambda_{\mathrm{tot}}-n}{\lambda_{\mathrm{tot}}}\in(0,1).
\]
Exponential Markov inequality and
\(e^{-\gamma_{\mathrm{sh}}}
\le1-\gamma_{\mathrm{sh}}+\gamma_{\mathrm{sh}}^2/2\) give
\[
\Pr\{\operatorname{Pois}(\lambda_{\mathrm{tot}})<n\}
\le
\exp\left\{-\frac{(\lambda_{\mathrm{tot}}-n)^2}
                       {2\lambda_{\mathrm{tot}}}\right\}
\le
\exp\left\{-\frac{n(B-2)^2}{4B}\right\}.
\]
The last inequality uses that
\((\lambda_{\mathrm{tot}}-n)^2/(2\lambda_{\mathrm{tot}})\)
is increasing for \(\lambda_{\mathrm{tot}}>n\).

Define the shortage penalty by
\[
p_{\mathrm{sh}}(n,B)
=\exp\left\{-\frac{n(B-2)^2}{4B}\right\}.
\]
Integrating the reconstruction bound under either supported prior and then taking the fixed-sample estimator infimum proves
\[
c_{\mathrm{sp}}\frac M{K_{\mathrm{sp}}^2}
-p_{\mathrm{sh}}(n,B)
\le\mathfrak R_{n,d,\epsilon}
\]
under the sparse conditions, and
\[
\frac5{2560000}\frac{\mathfrak a^2}{K_{\mathrm{de}}^2}
-p_{\mathrm{sh}}(n,B)
\le\mathfrak R_{n,d,\epsilon}
\]
under the dense conditions.

For the sparse bound choose
\[
M_{\mathrm{sp}}
=\min\left\{
K_{\mathrm{sp}}^2,
\frac{\eta_{\mathrm{sp}}dK_{\mathrm{sp}}}{Bn\epsilon}
\right\}.
\]
Whenever the second entry is at least two, this is an admissible support endpoint, and
\[
c_{\mathrm{sp}}
\min\left\{1,\frac{\eta_{\mathrm{sp}}d}
                    {Bn\epsilon K_{\mathrm{sp}}}\right\}
-p_{\mathrm{sh}}(n,B)
\le\mathfrak R_{n,d,\epsilon}.
\]
For the dense bound choose
\[
\mathfrak a_{\mathrm{de}}^2
=\min\left\{\frac14,
\frac{\eta_{\mathrm{de}}dK_{\mathrm{de}}}{Bn\epsilon}\right\}.
\]
Then \(0<\mathfrak a_{\mathrm{de}}\le1/2\) and the likelihood condition holds. Dividing by \(K_{\mathrm{de}}^2\) gives
\[
\frac5{2560000}
\min\left\{
\frac1{4K_{\mathrm{de}}^2},
\frac{\eta_{\mathrm{de}}d}{Bn\epsilon K_{\mathrm{de}}}
\right\}
-p_{\mathrm{sh}}(n,B)
\le\mathfrak R_{n,d,\epsilon}.
\]
% lean: CausalSmith.Stat.OptvalueVanishingoverlapRate.sparse_fixedSample_calibrated_lower

\item
\textbf{Fixed-sample transfer.}\quad
We join these two ranges while retaining their degree comparison. Define
\[
D_{\mathrm{lg}}=\max\{D_{\mathrm{sp}},D_{\mathrm{de}}\},
\qquad
q_{\mathrm{deg}}=\frac{C_{\mathrm{de}}+2}{C_{\mathrm{sp}}},
\]
\[
b_{\mathrm{splice}}
=\min\left\{
\frac{\eta_{\mathrm{sp}}}{8q_{\mathrm{deg}}^2},
\frac{\eta_{\mathrm{de}}}{q_{\mathrm{deg}}}
\right\},
\qquad
c_{\mathrm{splice}}
=\min\left\{
c_{\mathrm{sp}}\min\{1,\eta_{\mathrm{sp}}\},
\frac5{2560000}b_{\mathrm{splice}}
\right\}>0,
\]
\[
x_{\mathrm{splice}}=\frac d{Bn\epsilon}.
\]
The degree bounds imply
\[
K_{\mathrm{de}}\le q_{\mathrm{deg}}K_{\mathrm{sp}}.
\]

If \(\eta_{\mathrm{sp}}x_{\mathrm{splice}}K_{\mathrm{sp}}\ge2\), use the sparse bound. The elementary inequality
\[
\min\{1,\eta_{\mathrm{sp}}\}
\min\left\{1,\frac{x_{\mathrm{splice}}}{K_{\mathrm{sp}}}\right\}
\le
\min\left\{1,
\frac{\eta_{\mathrm{sp}}x_{\mathrm{splice}}}{K_{\mathrm{sp}}}\right\}
\]
shows that its leading term is at least
\(c_{\mathrm{splice}}\min\{1,x_{\mathrm{splice}}/K_{\mathrm{sp}}\}\).

Otherwise,
\(\eta_{\mathrm{sp}}x_{\mathrm{splice}}K_{\mathrm{sp}}\le2\).
The degree comparison and the definition of \(b_{\mathrm{splice}}\) give
\[
b_{\mathrm{splice}}
\frac{x_{\mathrm{splice}}}{K_{\mathrm{sp}}}
\le\frac1{4K_{\mathrm{de}}^2},
\qquad
b_{\mathrm{splice}}
\frac{x_{\mathrm{splice}}}{K_{\mathrm{sp}}}
\le\frac{\eta_{\mathrm{de}}x_{\mathrm{splice}}}{K_{\mathrm{de}}}.
\]
For the first inequality, multiply by \(4K_{\mathrm{de}}^2\) and use
\[
\frac{\eta_{\mathrm{sp}}}{8q_{\mathrm{deg}}^2}
\frac{x_{\mathrm{splice}}}{K_{\mathrm{sp}}}
\,4K_{\mathrm{de}}^2
\le\frac{\eta_{\mathrm{sp}}x_{\mathrm{splice}}K_{\mathrm{sp}}}{2}
\le1.
\]
For the second, use
\(K_{\mathrm{de}}\le q_{\mathrm{deg}}K_{\mathrm{sp}}\).
The dense bound therefore supplies the same leading term. In both cases,
\[
c_{\mathrm{splice}}
\min\left\{1,\frac d{Bn\epsilon K_{\mathrm{sp}}}\right\}
-p_{\mathrm{sh}}(n,B)
\le\mathfrak R_{n,d,\epsilon}
\qquad(d\ge D_{\mathrm{lg}}).
\]

Set
\[
c_{\mathrm{sh}}=\frac{c_{\mathrm{splice}}}{C_{\mathrm{sp}}+2}.
\]
Since \(K_{\mathrm{sp}}\le(C_{\mathrm{sp}}+2)L_d\) and
\(1/\{B(C_{\mathrm{sp}}+2)\}\le1\),
\[
\frac1{B(C_{\mathrm{sp}}+2)}\rho_{n,d,\epsilon}
\le
\min\left\{1,\frac d{Bn\epsilon K_{\mathrm{sp}}}\right\}.
\]
Thus
\[
\frac{c_{\mathrm{sh}}}{B}\rho_{n,d,\epsilon}
-p_{\mathrm{sh}}(n,B)
\le\mathfrak R_{n,d,\epsilon}
\qquad(d\ge D_{\mathrm{lg}},\ B>2).
\]
% lean: CausalSmith.Stat.OptvalueVanishingoverlapRate.largeAlphabet_lower_with_shortage

\item
Choose once and for all
\[
B_{\mathrm{lg}}=\max\{4,2048/c_{\mathrm{sh}}\},
\qquad
c_{\mathrm{lg}}=\frac{c_{\mathrm{sh}}}{2B_{\mathrm{lg}}}>0.
\]
For \(B\ge4\), define
\[
a_{\mathrm{sh}}=\frac{(B-2)^2}{8B}\ge\frac B{32}.
\]
The inequality \(e^{-v}\le1/v\) for \(v>0\) gives
\[
e^{-a_{\mathrm{sh}}n}\le\frac{32}{Bn},
\qquad
e^{-a_{\mathrm{sh}}n}\le e^{-a_{\mathrm{sh}}}\le\frac{32}{B}.
\]
Multiplying these bounds proves
\[
p_{\mathrm{sh}}(n,B)
=e^{-a_{\mathrm{sh}}n}e^{-a_{\mathrm{sh}}n}
\le\frac{1024}{B^2n}.
\]
Also \(\log d\le d-1\), so \(L_d\le d\). Since \(n\ge1\) and
\(\epsilon\le1/2\le1\),
\[
\frac1n\le1,
\qquad
\frac1n\le\frac d{n\epsilon L_d},
\qquad
\frac1n\le\rho_{n,d,\epsilon}.
\]
Because \(B_{\mathrm{lg}}c_{\mathrm{sh}}\ge2048\),
\[
p_{\mathrm{sh}}(n,B_{\mathrm{lg}})
\le\frac{1024}{B_{\mathrm{lg}}^2n}
\le\frac{c_{\mathrm{sh}}}{2B_{\mathrm{lg}}}\frac1n
\le c_{\mathrm{lg}}\rho_{n,d,\epsilon}.
\]
The preceding lower bound has leading coefficient
\(c_{\mathrm{sh}}/B_{\mathrm{lg}}=2c_{\mathrm{lg}}\). Subtracting the shortage penalty therefore gives
\[
c_{\mathrm{lg}}\rho_{n,d,\epsilon}
\le\mathfrak R_{n,d,\epsilon}
\qquad(d\ge D_{\mathrm{lg}}).
\]
% lean: CausalSmith.Stat.OptvalueVanishingoverlapRate.largeAlphabet_minimaxRisk_lower

\item
For the bounded-alphabet comparison, define
\[
r_{\mathrm{par}}=\min\left\{1,\frac1{n\epsilon}\right\},
\qquad
a_{\mathrm{par}}=\sqrt{\frac{r_{\mathrm{par}}}{64}}.
\]
Then
\[
0<a_{\mathrm{par}}\le\frac12,
\qquad
n\epsilon a_{\mathrm{par}}^2\le\frac1{16}.
\]
Use two dense observed laws \(\mathbb P_{\mathrm{par},1}\) and
\(\mathbb P_{\mathrm{par},0}\) with uniform covariate probabilities, propensity
\(\epsilon\), control mean \(1/2\), and treated means respectively
\(1/2+a_{\mathrm{par}}\) and \(1/2\). Define their scalar targets by
\[
\theta_{\mathrm{par},1}
=\Psi(\mathbb P_{\mathrm{par},1})=\frac12+a_{\mathrm{par}},
\qquad
\theta_{\mathrm{par},0}
=\Psi(\mathbb P_{\mathrm{par},0})=\frac12.
\]
For every \(x\in[d]\), the atom probabilities of these laws are
\[
\begin{aligned}
\mathbb P_{\mathrm{par},1}(x,0,0)
&=\mathbb P_{\mathrm{par},1}(x,0,1)
=\mathbb P_{\mathrm{par},0}(x,0,0)
=\mathbb P_{\mathrm{par},0}(x,0,1)
=\frac{1-\epsilon}{2d},\\
\mathbb P_{\mathrm{par},0}(x,1,0)
&=\mathbb P_{\mathrm{par},0}(x,1,1)=\frac{\epsilon}{2d},\\
\mathbb P_{\mathrm{par},1}(x,1,0)
&=\frac{\epsilon}{d}\left(\frac12-a_{\mathrm{par}}\right),
\qquad
\mathbb P_{\mathrm{par},1}(x,1,1)
=\frac{\epsilon}{d}\left(\frac12+a_{\mathrm{par}}\right).
\end{aligned}
\]
All reference atom probabilities are positive. Each untreated atom has zero probability difference and hence zero chi-square contribution. Each treated atom contributes
\[
\frac{(\epsilon a_{\mathrm{par}}/d)^2}{\epsilon/(2d)}
=\frac{2\epsilon a_{\mathrm{par}}^2}{d}.
\]
Summing over both treated outcomes and all \(d\) cells gives the finite-atom calculation
\[
\begin{aligned}
\chi^2(\mathbb P_{\mathrm{par},1},\mathbb P_{\mathrm{par},0})
&=\sum_{x\in[d]}\sum_{a,y\in\{0,1\}}
\frac{\{\mathbb P_{\mathrm{par},1}(x,a,y)
-\mathbb P_{\mathrm{par},0}(x,a,y)\}^2}
 {\mathbb P_{\mathrm{par},0}(x,a,y)}\\
&=d\cdot2\cdot\frac{2\epsilon a_{\mathrm{par}}^2}{d}
=4\epsilon a_{\mathrm{par}}^2.
\end{aligned}
\]
Define the two \(n\)-observation laws by
\[
\mathbb Q_{\mathrm{par},\iota}
=\mathbb P_{\mathrm{par},\iota}^{\otimes n},
\qquad \iota\in\{0,1\}.
\]
Product likelihoods give
\[
1+\chi^2(\mathbb Q_{\mathrm{par},1},\mathbb Q_{\mathrm{par},0})
=(1+4\epsilon a_{\mathrm{par}}^2)^n
\le e^{4n\epsilon a_{\mathrm{par}}^2}
\le e^{1/4}\le\frac43.
\]
In particular their chi-square divergence is at most one, so their total variation is at most \(1/2\). Thresholding an arbitrary estimator at
\(1/2+a_{\mathrm{par}}/2\) shows that one of the two laws has probability at least \(1/4\) of absolute error at least \(a_{\mathrm{par}}/2\). Hence
\[
\max_{\iota=0,1}
E_{\mathbb Q_{\mathrm{par},\iota}}
\{\widehat V-\theta_{\mathrm{par},\iota}\}^2
\ge\frac{a_{\mathrm{par}}^2}{16}
=\frac1{1024}r_{\mathrm{par}}.
\]
Taking the estimator infimum gives
\[
\frac1{1024}\min\left\{1,\frac1{n\epsilon}\right\}
\le\mathfrak R_{n,d,\epsilon}
\qquad(d\ge2).
\]

Now define
\[
c_{\mathrm{all}}
=\min\left\{c_{\mathrm{lg}},\frac1{1024D_{\mathrm{lg}}}\right\}>0.
\]
For \(d\ge D_{\mathrm{lg}}\), the large-alphabet bound applies.
For \(2\le d<D_{\mathrm{lg}}\), using \(L_d\ge1\) gives
\[
\rho_{n,d,\epsilon}
\le
D_{\mathrm{lg}}\min\left\{1,\frac1{n\epsilon}\right\}.
\]
Indeed, \(\rho_{n,d,\epsilon}\le1\le D_{\mathrm{lg}}\), and
\[
\rho_{n,d,\epsilon}
\le\frac d{n\epsilon L_d}
\le\frac{D_{\mathrm{lg}}}{n\epsilon}.
\]
Thus in this second case,
\[
c_{\mathrm{all}}\rho_{n,d,\epsilon}
\le
\frac1{1024}\min\left\{1,\frac1{n\epsilon}\right\}
\le\mathfrak R_{n,d,\epsilon}.
\]
Together the two cases establish
\[
c_{\mathrm{all}}\rho_{n,d,\epsilon}
\le\mathfrak R_{n,d,\epsilon}
\qquad(n\ge1,\ d\ge2,\ 0<\epsilon\le1/2).
\]
% lean: CausalSmith.Stat.OptvalueVanishingoverlapRate.allAlphabet_minimaxRisk_lower

\item
\textbf{Final consequences.}\quad
By \cref{obj:thm:observable-upper}, there are universal calibrations
\(H_0>0,\kappa>0,D_0\ge2\) and a universal \(C_{\mathrm u}>0\) for the measurable estimator in
\cref{obj:def:armwise-estimator}. Define its worst-case observed risk by
\[
\mathcal R^{\mathrm{AF}}_{n,d,\epsilon}
=
\sup_{\mathbb P\in\mathcal D_{d,\epsilon}^{\mathrm{obs}}}
E_{\mathbb P}
\left[
\{\widehat V_{n,d,\epsilon}^{\mathrm{AF}}-\Psi(\mathbb P)\}^2
\right].
\]
Apply the cited bound to the product sampling law
\(\mathbb P^{\otimes n}\), which satisfies
\cref{obj:ass:iid-sampling}. This gives
\[
\mathcal R^{\mathrm{AF}}_{n,d,\epsilon}
\le C_{\mathrm u}\rho_{n,d,\epsilon}.
\]
Set
\[
c_\star=\min\{c_{\mathrm{all}},C_{\mathrm u}\},
\qquad
C_\star=C_{\mathrm u}.
\]
Then \(0<c_\star\le C_\star<\infty\). Nonnegativity of
\(\rho_{n,d,\epsilon}\) permits reducing the lower constant, and the estimator infimum in
\cref{obj:def:minimax-risk} is bounded by the risk of this particular estimator. Consequently,
\[
c_\star\rho_{n,d,\epsilon}
\le\mathfrak R_{n,d,\epsilon}
\le\mathcal R^{\mathrm{AF}}_{n,d,\epsilon}
\le C_\star\rho_{n,d,\epsilon}.
\]

For each measurable \(\widehat V\), identification and completion in
\cref{obj:prop:identification} give the equality
\[
\begin{aligned}
&\sup_{P\in\mathcal P_{d,\epsilon}}
E_{\operatorname{Obs}(P)}
[\{\widehat V-V^\star(P)\}^2]\\
&\hspace{2cm}=
\sup_{\mathbb P\in\mathcal D_{d,\epsilon}^{\mathrm{obs}}}
E_{\mathbb P}
[\{\widehat V-\Psi(\mathbb P)\}^2].
\end{aligned}
\]
For one inequality use
\(V^\star(P)=\Psi\{\operatorname{Obs}(P)\}\) and the observed marginal's membership in the observed class. For the other, complete each observed law to a causal law with that same marginal and value. Taking the estimator infimum preserves equality, so the observed bounds also prove the two stated causal minimax bounds.
% lean: CausalSmith.Stat.OptvalueVanishingoverlapRate.weighted_separation_and_frontier

\item
Finally, let the sequences in the statement be given. Define
\[
r_k=\frac{d_k}{n_k\epsilon_k\log(ed_k)},
\qquad
\rho_k=\min\{1,r_k\}.
\]
For the finitely many possible indices with \(n_k=0\), use the totalized division convention \(r_k=0\). Since \(n_k\to\infty\), one has \(n_k\ge1\) eventually. For all \(k\), \(r_k,\rho_k\ge0\), and
\[
\rho_k\longrightarrow0
\quad\Longleftrightarrow\quad
r_k\longrightarrow0.
\]
The reverse implication follows from \(0\le\rho_k\le r_k\).
For the forward implication, eventually \(\rho_k<1\); then
\(r_k\ge1\) would force \(\rho_k=1\), so eventually
\(r_k<1\) and \(\rho_k=r_k\).

Suppose a measurable estimator sequence \((\widehat V_k)\) has worst-case risks
\[
\mathcal W_k
=
\sup_{\mathbb P\in\mathcal D_{d_k,\epsilon_k}^{\mathrm{obs}}}
E_{\mathbb P}
[\{\widehat V_k-\Psi(\mathbb P)\}^2]
\longrightarrow0.
\]
The sample alphabet is finite. Thus each \(\widehat V_k\) is bounded on it, and these squared-error suprema are finite and nonnegative. The lower bound and the estimator infimum give, eventually,
\[
0\le\rho_k
\le\frac{\mathfrak R_{n_k,d_k,\epsilon_k}}{c_\star}
\le\frac{\mathcal W_k}{c_\star}.
\]
The squeeze theorem implies \(\rho_k\to0\), hence \(r_k\to0\).

Conversely, suppose \(r_k\to0\). Define a measurable estimator for every index by
\[
\widehat V_k=
\begin{cases}
\widehat V_{n_k,d_k,\epsilon_k}^{\mathrm{AF}},&n_k\ge1,\\
0,&n_k=0\text{ and }d_k<D_0,\\
1/2,&n_k=0\text{ and }d_k\ge D_0.
\end{cases}
\]
The zero-sample branches define constants on the singleton empty-sample space. For the eventual positive-sample indices, the established upper bound gives
\[
0\le
\sup_{\mathbb P\in\mathcal D_{d_k,\epsilon_k}^{\mathrm{obs}}}
E_{\mathbb P}
[\{\widehat V_k-\Psi(\mathbb P)\}^2]
\le C_\star\rho_k\longrightarrow0.
\]
A final application of the squeeze theorem proves uniform consistency and completes the equivalence.
% lean: CausalSmith.Stat.OptvalueVanishingoverlapRate.uniformConsistency_iff_of_frontier
\end{enumerate}
\end{proof}

\begin{proof}[Proof of \cref{obj:thm:all-estimator-lower}]
Choose the universal lower comparison constant \(c_\star>0\) supplied by the observed-risk conclusion of \cref{obj:thm:weighted-separation-and-frontier}. For arbitrary \(n,d\in\mathbb N\) and \(\epsilon\in\mathbb R\) satisfying \(n\ge1\), \(d\ge2\), and \(0<\epsilon\le1/2\), that conclusion gives
\[
c_\star\min\left\{1,\frac{d}{n\epsilon\log(ed)}\right\}
\le \mathfrak R_{n,d,\epsilon}.
\]
% lean: CausalSmith.Stat.OptvalueVanishingoverlapRate.weighted_separation_and_frontier
Since \(L_d=\log(ed)\), this is the required inequality. The constant \(c_\star\) is positive and independent of \(n,d,\epsilon\), as required.
% lean: CausalSmith.Stat.OptvalueVanishingoverlapRate.all_estimator_lower
\end{proof}
